% Options for packages loaded elsewhere
\PassOptionsToPackage{unicode}{hyperref}
\PassOptionsToPackage{hyphens}{url}
%
\documentclass[
  10pt,
]{article}
\usepackage{lmodern}
\usepackage{amssymb,amsmath}
\usepackage{ifxetex,ifluatex}
\ifnum 0\ifxetex 1\fi\ifluatex 1\fi=0 % if pdftex
  \usepackage[T1]{fontenc}
  \usepackage[utf8]{inputenc}
  \usepackage{textcomp} % provide euro and other symbols
\else % if luatex or xetex
  \usepackage{unicode-math}
  \defaultfontfeatures{Scale=MatchLowercase}
  \defaultfontfeatures[\rmfamily]{Ligatures=TeX,Scale=1}
\fi
% Use upquote if available, for straight quotes in verbatim environments
\IfFileExists{upquote.sty}{\usepackage{upquote}}{}
\IfFileExists{microtype.sty}{% use microtype if available
  \usepackage[]{microtype}
  \UseMicrotypeSet[protrusion]{basicmath} % disable protrusion for tt fonts
}{}
\makeatletter
\@ifundefined{KOMAClassName}{% if non-KOMA class
  \IfFileExists{parskip.sty}{%
    \usepackage{parskip}
  }{% else
    \setlength{\parindent}{0pt}
    \setlength{\parskip}{6pt plus 2pt minus 1pt}}
}{% if KOMA class
  \KOMAoptions{parskip=half}}
\makeatother
\usepackage{xcolor}
\IfFileExists{xurl.sty}{\usepackage{xurl}}{} % add URL line breaks if available
\IfFileExists{bookmark.sty}{\usepackage{bookmark}}{\usepackage{hyperref}}
\hypersetup{
  hidelinks,
  pdfcreator={LaTeX via pandoc}}
\urlstyle{same} % disable monospaced font for URLs
\usepackage[a4paper,margin=18mm]{geometry}
\setlength{\emergencystretch}{3em} % prevent overfull lines
\providecommand{\tightlist}{%
  \setlength{\itemsep}{0pt}\setlength{\parskip}{0pt}}
\setcounter{secnumdepth}{-\maxdimen} % remove section numbering

\usepackage{esint}
\usepackage{xeCJK}
\setmainfont{DejaVu Serif}
\setmonofont{DejaVu Sans Mono}
\setCJKmainfont{Noto Serif CJK SC}
\setCJKmonofont{Noto Sans Mono CJK SC}
\setmathfont{Asana Math}
\setmathfont{Latin Modern Math}[range=\lrcorner]
\setlength{\emergencystretch}{3em}
\allowdisplaybreaks
\hypersetup{hypertexnames=false}

\author{}
\date{}

\begin{document}

\hypertarget{technical-supplement-intrinsic-calibration-and-parameter-theory}{%
\section{Technical supplement: intrinsic calibration and parameter
theory}\label{technical-supplement-intrinsic-calibration-and-parameter-theory}}

Research snapshot: 844e62f7304d9f890914a6a4370ec39291ddc26e.

Selected complete older proofs needed by the recent papers, not newly
counted results. Historical endpoint limitations refer to their date and
are updated by Volumes 1-3.

Read 00\_READ\_FIRST.md for superseded statements. Each source below is
complete and byte-preserved. Original local links remain relative to
their source; locate other texts using Volume 6. Formula numbers restart
in each source.

\begin{center}\rule{0.5\linewidth}{0.5pt}\end{center}

\hypertarget{source-1-line1_integral_rigidity_paper.md}{%
\subsection{Source 1:
line1\_integral\_rigidity\_paper.md}\label{source-1-line1_integral_rigidity_paper.md}}

Original path: mathematics/current/line1\_integral\_rigidity\_paper.md

SHA-256:
ca664aa1f07d2907c8bd0b346d4e0b824cb858335fa2a8d15c9cce33ba7e3d54

BEGIN COMPLETE SOURCE

\hypertarget{integral-einstein-rigidity-from-polarized-monopole-classes}{%
\section{Integral Einstein rigidity from polarized monopole
classes}\label{integral-einstein-rigidity-from-polarized-monopole-classes}}

\hypertarget{abstract}{%
\subsection{Abstract}\label{abstract}}

We prove a scale-invariant, integral obstruction to Einstein metrics on
closed oriented four-manifolds with suitable polarized Seiberg--Witten
classes. The analytic statement uses actual solutions for every positive
weight, but requires no Sobolev, injectivity-radius, diameter, volume,
pointwise-curvature or curvature-derivative bound. Its proof retains the
nonnegative errors in the monopole inequalities, bounds the spinor by a
truncated weight, constructs a globally smooth polynomial Weyl model,
and tests the contracted Bianchi identity against the divergence of that
model. A bundle Hessian estimate proved below closes this self-test
without a geometric compactness argument.

The resulting universal constant gives an integral Einstein gap for the
connected sums occurring in LeBrun's Einstein obstruction. The class may
depend on the metric, including its chamber when the positive Betti
number is one; the same selected class must work for every weight in
that metric's conformal class. Positive and negative Einstein constants
and the separate Ricci-flat case are treated explicitly. The
positive-budget application strengthens the numerical obstruction in
LeBrun's Theorem 3.3 and Corollary 3.4. Negative topological budgets are
handled directly by Chern--Gauss--Bonnet; the zero-budget, zero-surgery
corner is not an Einstein nonexistence statement.

This manuscript makes the inherited integral-rigidity argument
self-contained analytically. It does not assert priority for either that
argument or the classical gauge-theoretic inputs.

\hypertarget{conventions-and-statements}{%
\subsection{1. Conventions and
statements}\label{conventions-and-statements}}

Let \(M\) be a closed, connected, oriented smooth four-manifold, and let

\[
 V=\operatorname{Vol}_g(M),\qquad T=2\chi(M)+3\tau(M),\qquad
 P=\int_M|\operatorname{Ric}^{\circ}|^2.
\]

Integrals and norms use \(g\), unless stated otherwise. Ricci is a
covariant two-tensor; \(W^+\) is the trace-free symmetric curvature
endomorphism of \(\Lambda^+\), with its operator Hilbert--Schmidt norm.
The Laplacian is nonnegative. Write \(x_-:=\max\{-x,0\}\) and

\[
 4\pi^2T=\int_M\left(\frac{s^2}{24}+2|W^+|^2
                          -\frac{|\operatorname{Ric}^{\circ}|^2}{2}\right).
 \tag{1.1}
\]

For a Spin\(^c\) structure with determinant line \(L\), put

\[
 a=c_1(L)\in H^2(M;\mathbb R),\qquad
 a^+=\operatorname{pr}_{H_g^+}a,\qquad p=(a^+)^2.
\]

All cohomological squares mean intersection squares. The quadratic
spinor map is normalized by

\[
 |\sigma(\Phi)|^2=|\Phi|^4/8.
\]

\textbf{Definition 1.1 (the conformal monopole hypothesis).} A pair
\((g,a)\) satisfies this hypothesis if \(p>0\) and one fixed Spin\(^c\)
structure with first Chern class \(a\) has an unperturbed
Seiberg--Witten solution for every metric conformal to \(g\).
Equivalently for our purposes, for every smooth \(f>0\) there is a
solution of

\[
 D_A\Phi=0,\qquad -iF_A^+=f\sigma(\Phi).                 \tag{1.2}
\]

The implication from conformal solvability to (1.2) is proved in §2. The
class is not allowed to change as \(f\) changes. A different class may
be selected for a different initial metric.

Define the scale-invariant integral Einstein error by

\[
 \mathcal E(g):=\inf_{\lambda\ne0}
  \frac{\int_M|\operatorname{Ric}_g-\lambda g|^2}
       {\lambda^2V}.                                      \tag{1.3}
\]

\textbf{Theorem 1.2 (integral rigidity).} There is an absolute constant

\[
                       \varepsilon_0>0                   \tag{1.4}
\]

with the following property. Suppose that for a metric \(g\) there is a
class \(a\) satisfying Definition 1.1 and

\[
                     (a^+)^2\ge A>0,\qquad a^2\ge T.
 \tag{1.5}
\]

Then

\[
 \boxed{\displaystyle
 \mathcal E(g)+\left(\frac TA-\frac23\right)_+
                         \ge\varepsilon_0.}              \tag{1.6}
\]

In particular, the conclusion holds for every metric if such a class can
be selected metricwise, with the same number \(A\). No regularity of the
class selection is needed. The constant is existential and universal;
this paper does not assign it a numerical value.

The main analytic ingredient is a calibration gap. Let

\[
 \begin{aligned}
 \mathcal O&=\operatorname{Orb}\operatorname{diag}
                    (-1/6,1/12,1/12)\subset S^2_0(\mathbb R^3),\\
 F&=\operatorname{dist}_{\mathrm{HS}}(W^+,\mathcal O)^2,\qquad
 E=\operatorname{Ric}+g,\qquad e=s+4,\\
 \alpha&=\sqrt{2/3}|W^+|,\qquad q=-s+3\alpha,\\
 \mathcal C&=\int_M(|E|^2+F),\qquad
 \Delta_+=\int_Mq_+^2-72\pi^2p,\\
 \kappa&=(\mathcal C+\Delta_+)/V,\qquad
 B_a=(-a^2)_+/V.
 \end{aligned}                                              \tag{1.7}
\]

Here \(\mathcal O\) is interpreted in each oriented orthonormal frame of
\(\Lambda^+\); the distance is frame independent.

\textbf{Theorem 1.3 (analytic calibration gap).} Under Definition 1.1,

\[
 \Delta_+\ge0,\qquad (1+B_a)\kappa\ge c_{\mathrm{cal}}>0, \tag{1.8}
\]

where \(c_{\mathrm{cal}}\) is absolute.

We prove this theorem first. Throughout the analytic proof, \(C\)
denotes an absolute positive constant whose value can change from line
to line. Explicit intermediate constants are retained where they clarify
the topological transfer, not to advertise an optimized final constant.

\hypertarget{actual-weighted-solutions-and-the-retained-monopole-identities}{%
\subsection{2. Actual weighted solutions and the retained monopole
identities}\label{actual-weighted-solutions-and-the-retained-monopole-identities}}

\hypertarget{why-the-weighted-equation-has-solutions}{%
\subsubsection{2.1. Why the weighted equation has
solutions}\label{why-the-weighted-equation-has-solutions}}

Let \(\widehat g=e^{2u}g=f^{-2}g\), so \(u=-\log f\). Under the
fiberwise unitary spinor identification \(\beta\), the conformal
formulas are

\[
 \begin{aligned}
 D_A^{\widehat g}(e^{-3u/2}\beta\Phi)
        &=e^{-5u/2}\beta D_A^g\Phi,\\
 \sigma_{\widehat g}(\beta v)&=e^{2u}\sigma_g(v)
                   \quad\hbox{as covariant two-forms}.
 \end{aligned}                                              \tag{2.1}
\]

For an ordinary solution \((A,\widehat\Phi)\) for \(\widehat g\), set

\[
                 \Phi=f^{-3/2}\beta^{-1}\widehat\Phi.
\]

Then \(D_A^g\Phi=0\), and

\[
 \sigma_g(\Phi)=f^{-1}\sigma_{\widehat g}(\widehat\Phi).
\]

The Hodge star on two-forms is conformally invariant in dimension four,
so the ordinary curvature equation becomes (1.2). In particular
\(H_g^+\) does not change during this operation. This is the conformal
weighted monopole construction of
\href{https://arxiv.org/pdf/math/0611450}{LeBrun, Lemma 3.7}; the
displayed formulas specify the normalization used here.

Set

\[
             \rho=|\Phi|^2,\qquad \psi=2\sqrt2\sigma(\Phi).
 \tag{2.2}
\]

Thus \(\psi\) is a smooth self-dual two-form, \(|\psi|=\rho\), and
\(\rho\) itself is smooth, including at spinor zeros. The harmonic
projection of the determinant curvature gives

\[
 \int_M f^2\rho^2=8\|F_A^+\|_2^2\ge32\pi^2p.           \tag{2.3}
\]

Only the squared norm is used, so the equivalent opposite sign
convention for the Chern representative does not change this formula.

\hypertarget{nonnegative-errors}{%
\subsubsection{2.2. Nonnegative errors}\label{nonnegative-errors}}

The spinor Weitzenböck equation and the self-dual Hodge formula are

\[
 4\nabla_A^*\nabla_A\Phi+(s+f\rho)\Phi=0,                \tag{2.4}
\]

\[
 \|d\psi\|_2^2+\|d^*\psi\|_2^2
   =\int_M\left(|\nabla\psi|^2+\frac{s}{3}\rho^2
                                  -2\langle W^+\psi,\psi\rangle\right).
 \tag{2.5}
\]

The differential of the Hopf quadratic map in (2.2) satisfies

\[
 |d\psi_z(v)|^2
       =4|z|^2|v|^2-4\bigl(\operatorname{Im}\langle z,v\rangle\bigr)^2
       \le4|z|^2|v|^2.                                   \tag{2.6}
\]

The same identity applies covariantly because the map is equivariant and
parallel. Pairing (2.4) with \(\rho\Phi\) and integrating yields

\[
 2\int_M|d\rho|^2+4\int_M\rho|\nabla_A\Phi|^2
                  +\int_Ms\rho^2+\int_M f\rho^3=0.       \tag{2.7}
\]

Let \(w\) be the smallest eigenvalue of \(W^+\). For a trace-free
symmetric three-by-three matrix, \(w\ge-\alpha\). Define

\[
 \begin{aligned}
 D_{\mathrm{amp}}&=2\int_M|d\rho|^2,\\
 D_{\mathrm{quad}}&=\int_M(4\rho|\nabla_A\Phi|^2-|\nabla\psi|^2),\\
 D_H&=\|d\psi\|_2^2+\|d^*\psi\|_2^2=2\|d\psi\|_2^2,\\
 D_{\mathrm{eig}}&=2\int_M
           (\langle W^+\psi,\psi\rangle+\alpha\rho^2),\\
 D_0&=D_{\mathrm{amp}}+D_{\mathrm{quad}}+D_H+D_{\mathrm{eig}}.
 \end{aligned}                                              \tag{2.8}
\]

Each term is nonnegative. Equations (2.5)--(2.7) give the exact identity

\[
 D_0=\int_M\left(\frac23q\rho^2-f\rho^3\right).
 \tag{2.9}
\]

For \(v=3\rho/2\), this is

\[
                  \int_M(qv^2-fv^3)=\frac{27}{8}D_0.
 \tag{2.10}
\]

\hypertarget{signed-spectral-and-scalar-estimates}{%
\subsubsection{2.3. Signed spectral and scalar
estimates}\label{signed-spectral-and-scalar-estimates}}

Put

\[
                  b=2s/3+2w,\qquad h=b_-.
\]

Using \(w\), rather than \(-\alpha\), in (2.5)--(2.7) gives

\[
                       \int_M f\rho^3\le\int_M h\rho^2.
 \tag{2.11}
\]

Choose positive smooth functions \(h_j\ge h\) converging uniformly to
\(h\), and use the actual solutions with \(f_j=\sqrt{h_j}\). Set

\[
 C_j=\int_M f_j^2\rho^2,\quad Z_j=\int_M f_j\rho^3,
 \quad Q_j=\int_M f_j^4.
\]

Equation (2.11) and Hölder imply

\[
 32\pi^2p\le C_j\le Q_j^{1/3}Z_j^{2/3}
                            \le Q_j^{1/3}C_j^{2/3}.
\]

Consequently \(32\pi^2p\le Q_j\), and uniform approximation proves

\[
 \int_M(2s/3+2w)_-^2\ge32\pi^2p,\qquad
 \int_M q_+^2\ge72\pi^2p.                               \tag{2.12}
\]

The second follows from \(b_-\le(2/3)q_+\). This argument passes to a
limit of inequalities, not to a limiting spinor or connection.

For \(f=1\), pairing (2.4) with \(\Phi\) gives

\[
 \int_M\rho^2\le\int_Ms_-\rho
        \le\|s_-\|_2\|\rho\|_2.
\]

Together with (2.3), this proves

\[
                        \int_Ms_-^2\ge32\pi^2p.          \tag{2.13}
\]

In particular \(p>0\) excludes a scalar-flat metric. These are the
classical monopole curvature estimates in the normalization needed here;
the retained terms in (2.8) will supply their quantitative use.

\hypertarget{a-bounded-spinor-with-small-retained-errors}{%
\subsection{3. A bounded spinor with small retained
errors}\label{a-bounded-spinor-with-small-retained-errors}}

The weights used in §2.3 need not bound the spinor uniformly. For the
calibration argument we instead set

\[
       Q=\max\{q,3\},\qquad
       f_0=\max\{\sqrt Q,s_-/4\}.                        \tag{3.1}
\]

Choose smooth majorants \(f_j\ge f_0\) with \(f_j\to f_0\) uniformly,
and choose an actual solution for each \(f_j\). At a positive maximum of
\(\rho\), (2.4) implies \(s+f_j\rho\le0\). Hence

\[
                           0\le\rho\le4,\qquad 0\le v\le6.
 \tag{3.2}
\]

All errors \(o_j(1)\) below tend to zero for this fixed smooth metric
and fixed class. They require only uniform approximation of \(f_0\), not
bounds on derivatives of \(f_j\). No assertion of convergence of the
solutions is made.

The two truncations cost only scalar error:

\[
 0\le Q^2-q_+^2\le9e^2,\qquad
 0\le f_0^2-Q\le\frac9{16}e^2.                          \tag{3.3}
\]

Indeed, the first difference vanishes unless \(q<3\), which implies
\(s>-3\) and \(e>1\). For the second, a nonzero difference implies
\(s_->4\sqrt3>6\), and then \(s_-\le3|s+4|\).

Define ideal moments using \(Q\), not the full truncated weight:

\[
 C_*:=\int_M Qv^2,\qquad Z_*:=\int_M\sqrt Qv^3,
 \qquad E_Q:=\int_M Q^2,\qquad P_*:=C_*-Z_*.
 \tag{3.4}
\]

Equation (2.10), \(Q\ge q\), and \(f_j\ge\sqrt Q\) give

\[
 P_*\ge\frac{27}{8}D_0\ge0.                            \tag{3.5}
\]

The actual curvature moment is

\[
 C_{\mathrm{act}}:=\int_M f_j^2v^2
                      =18\|F_A^+\|_2^2\ge72\pi^2p.
\]

By (3.2)--(3.3),

\[
 0\le C_{\mathrm{act}}-C_*
                     \le\frac{81}{4}\int_M e^2+o_j(1).
\]

Hölder gives \(C_*\le E_Q^{1/3}Z_*^{2/3}\). Since \(Z_*\le C_*\), it
follows that \(C_*\le E_Q\). Thus, with

\[
                        S_*:=\Delta_++\frac{117}{4}\int_Me^2,
\]

we have

\[
                     0\le E_Q-C_*\le S_*+o_j(1).         \tag{3.6}
\]

For completeness, two elementary moment inequalities recover all the
errors needed below. Write

\[
 \theta=(Z_*/E_Q)^{1/3},\qquad
 H_*=E_Q^{1/3}Z_*^{2/3}-C_*.
\]

Here \(0<\theta\le1\), \(H_*\ge0\), and direct algebra gives

\[
 E_Q-C_*=H_*+\frac{1+\theta}{\theta^2}(P_*+H_*)\ge2P_*.
 \tag{3.7}
\]

For the probability measure \(d\nu=Q^2d\mu/E_Q\), let \(r=v/\sqrt Q\),
\(c=\int r^2d\nu\), \(z=\int r^3d\nu\), and \(m=\int rd\nu\). We have
\(z\le c\le z^{2/3}\), hence \(c\le1\). Also \(c^2\le mz\le mc\), so
\(m\ge c\). Therefore

\[
 \int_M Q(v-\sqrt Q)^2=E_Q(c-2m+1)\le E_Q-C_*.
 \tag{3.8}
\]

Equations (3.5)--(3.8) prove

\[
 \begin{aligned}
 D_0&\le\frac4{27}S_*+o_j(1),\\
 \int_M|d\rho|^2,\ \int_M|d\psi|^2
     &\le\frac2{27}S_*+o_j(1),\\
 \int_MQ(v-\sqrt Q)^2&\le S_*+o_j(1).
 \end{aligned}                                              \tag{3.9}
\]

The point of the ideal moments is that these conclusions do not require
any estimate of \(\int f_j^4\), which could be unnecessarily large.

\hypertarget{a-globally-smooth-approximate-weyl-tensor}{%
\subsection{4. A globally smooth approximate Weyl
tensor}\label{a-globally-smooth-approximate-weyl-tensor}}

The norm of an element of \(\mathcal O\) is \(1/(2\sqrt6)\).
Consequently

\[
 \left\|q-\frac92\right\|_2^2
       \le2\|e\|_2^2+12\int_MF\le12\mathcal C.           \tag{4.1}
\]

Clipping at \(3\) moves \(q\) toward \(9/2\), so the same bound holds
with \(Q\). Since \(Q\ge3\), (3.8) implies

\[
 \|\rho-\sqrt2\|_2^2
  \le\frac8{27}S_*+\frac2{27}\left\|Q-\frac92\right\|_2^2+o_j(1).
 \tag{4.2}
\]

Here the first term follows by writing
\(\rho-\sqrt2=(2/3)(v-\sqrt Q)+(2/3)(\sqrt Q-\sqrt{9/2})\) and using
\(Q\ge3\).

To estimate direction, choose a measurable rank-one orthogonal
projection \(\Pi\) onto a lowest eigenspace of \(W^+\). Then

\[
                  A_0=I/12-\Pi/4
\]

is a nearest point of \(\mathcal O\). On the good set \(F\le1/256\),
comparison of ordered eigenvalues with \(A_0\) gives

\[
                   \lambda_2(W^+)-\lambda_1(W^+)\ge1/8.
\]

Putting \(y=(I-\Pi)\psi\), it follows from the eigenvalue error in (2.8)
that

\[
 \int_{\{F\le1/256\}}|y|^2\le\frac{16}{27}S_*+o_j(1).
\]

The complementary set has volume at most \(256\int F\), while
\(|\psi|^2\le16\). Thus

\[
                 \int_M|y|^2\le\frac{16}{27}S_*+4096\int_MF+o_j(1).
 \tag{4.3}
\]

No derivative of \(\Pi\) will be taken. Instead define the polynomial
endomorphism

\[
       A_\psi=\frac{\rho^2}{24}I-\frac18\psi\otimes\psi,
             \qquad R=W^+-A_\psi.                       \tag{4.4}
\]

It is globally smooth, trace-free, and defined even at spinor zeros. Its
eigenvalues are \(-\rho^2/12,\rho^2/24,\rho^2/24\), and
\(|A_\psi|^2=\rho^4/96\). The rank-one identity

\[
 |\psi\otimes\psi-\rho^2\Pi|^2=2\rho^2|y|^2
\]

gives the useful comparison

\[
 |A_\psi-A_0|^2
       \le\frac{\rho^2}{16}|y|^2+\frac{(\rho^2-2)^2}{48}.
 \tag{4.5}
\]

For example, (4.5) follows by inserting the intermediate matrix
\(\rho^2 A_0/2\) and applying \(|X+Y|^2\le2|X|^2+2|Y|^2\). Since
\(\rho\le4\) and \((4+\sqrt2)^2<36\), (4.1)--(4.5) imply

\[
 \begin{aligned}
 \|R\|_2^2
 &\le\frac{44}{27}S_*+8194\int_MF
                     +\frac19\left\|Q-\frac92\right\|_2^2+o_j(1)\\
 &\le2\Delta_++8500\mathcal C+o_j(1).
 \end{aligned}                                              \tag{4.6}
\]

In the final line we used \(\int e^2\le4\int|E|^2\). In particular, each
of

\[
 \|e\|_2^2,\quad\|\rho-\sqrt2\|_2^2,\quad\|R\|_2^2,
            \quad\|d\rho\|_2^2,\quad\|d\psi\|_2^2
\]

is bounded by \(C\kappa V+o_j(1)\).

\hypertarget{a-hessian-estimate-that-needs-no-sobolev-constant}{%
\subsection{5. A Hessian estimate that needs no Sobolev
constant}\label{a-hessian-estimate-that-needs-no-sobolev-constant}}

\textbf{Lemma 5.1.} Let \(u\) be a smooth section of a metric or
Hermitian vector bundle with compatible connection over a closed
Riemannian four-manifold. Let \(\mathscr F\) be its curvature, with
component norm \(\sum_{i,j}|\mathscr F_{ij}|_{\mathrm{HS}}^2\), and set
\(K=\|u\|_\infty\). Then

\[
 \begin{aligned}
 \|\nabla u\|_4^2&\le4K\|\nabla^2u\|_2,\\
 \|\nabla^2u\|_2^2
 &\le2\|\nabla^*\nabla u\|_2^2
       +32K^2\|\operatorname{Ric}\|_2^2
       +33K^2\|\mathscr F\|_2^2.
 \end{aligned}                                              \tag{5.1}
\]

\textbf{Proof.} Take real parts of Hermitian inner products when
necessary. Write \(v_i=\nabla_i u\), \(H_{ij}=\nabla_i v_j\), and
\(L=\sum_iH_{ii}=-\nabla^*\nabla u\). Integrating the divergence of
\(X_i=|v|^2\langle u,v_i\rangle\) gives

\[
 \int|v|^4
 \le K\int |v|^2(2|H|+|L|)
 \le K\|v\|_4^2(2\|H\|_2+\|L\|_2).
\]

Because \(|L|\le2|H|\) in dimension four, this proves the first
inequality, including the case \(\|v\|_4=0\).

For the second, antisymmetry gives

\[
 H_{ij}-H_{ji}=\mathscr F_{ij}u,\qquad
 \|H\|_2^2=\int\langle H_{ij},H_{ji}\rangle
                         +\frac12\sum_{i,j}\|\mathscr F_{ij}u\|_2^2.
 \tag{5.2}
\]

The commutator on the bundle-valued one-form \(v\) yields

\[
 \sum_i\nabla_i\nabla_jv_i
     =\nabla_jL+\operatorname{Ric}_{jk}v_k
                                      +\sum_i\mathscr F_{ij}v_i.
\]

Integrating by parts in the first term of (5.2) and using the
skew-adjointness of \(\mathscr F_{ij}\), we obtain the exact identity

\[
 \begin{aligned}
 \|H\|_2^2={}&\|L\|_2^2
  -\int\operatorname{Ric}_{ij}\langle v_i,v_j\rangle
  +\sum_{i,j}\int\langle\mathscr F_{ij}v_j,v_i\rangle
  +\frac12\sum_{i,j}\|\mathscr F_{ij}u\|_2^2.
 \end{aligned}                                              \tag{5.3}
\]

In particular no derivative of curvature appears. Hölder and the first
part bound the middle terms in absolute value by

\[
   4K(\|\operatorname{Ric}\|_2+\|\mathscr F\|_2)\|H\|_2.
\]

Use Young's inequality to absorb half of \(\|H\|_2^2\), and then
\((x+y)^2\le2(x^2+y^2)\). This gives exactly the second inequality in
(5.1). No Sobolev inequality or curvature smallness has been used.
\(\square\)

\hypertarget{applying-the-lemma-to-the-bounded-spinor}{%
\subsubsection{5.2. Applying the lemma to the bounded
spinor}\label{applying-the-lemma-to-the-bounded-spinor}}

The calibration budget gives

\[
 \begin{aligned}
 \int|W^+|^2&\le2\int F+V/12,\\
 \int|\operatorname{Ric}|^2&\le2\int|E|^2+8V,\\
 \int s^2&\le8\int|E|^2+32V.
 \end{aligned}                                              \tag{5.4}
\]

Equations (3.3) and (4.1) also imply

\[
                         \int f_j^2\le C(V+\mathcal C)+o_j(1).
 \tag{5.5}
\]

Indeed \(Q\le9/2+|Q-9/2|\); Cauchy--Schwarz and
\(\sqrt{V\mathcal C}\le(V+\mathcal C)/2\) suffice.

For the determinant connection, Chern--Weil and (3.2) give

\[
 \|F_A\|_2^2
   =2\|F_A^+\|_2^2-4\pi^2a^2
   \le4\int f_j^2+4\pi^2(-a^2)_+.
 \tag{5.6}
\]

The positive-spinor connection curvature \(\mathscr F_A\) is a fixed
linear representation of the Riemannian and determinant curvatures. In
the block decomposition of the Riemannian curvature operator, its
nontrivial spin action uses the \(\Lambda^+\) output: it sees \(W^+\),
\(s\), and the trace-free Ricci block, but not \(W^-\). Consequently
(5.4)--(5.6) imply

\[
             \int|\mathscr F_A|^2
               \le C\bigl(V+\mathcal C+(-a^2)_+\bigr)+o_j(1).
 \tag{5.7}
\]

From (2.4) and \(\rho\le4\),

\[
 \|\nabla_A^*\nabla_A\Phi\|_2^2
       \le\frac12\int s^2+8\int f_j^2
       \le C(V+\mathcal C)+o_j(1).                        \tag{5.8}
\]

Apply Lemma 5.1 with \(\|\Phi\|_\infty\le2\). Since \(A_\psi\) is a
parallel polynomial of degree four in \(\Phi\),

\[
 |\nabla^2A_\psi|
      \le C\bigl(|\Phi|^3|\nabla_A^2\Phi|
                               +|\Phi|^2|\nabla_A\Phi|^2\bigr).
\]

We conclude that

\[
 \int\bigl(|\nabla_A^2\Phi|^2+|\nabla_A\Phi|^4
                                      +|\nabla^2A_\psi|^2\bigr)
       \le C\bigl(V+\mathcal C+(-a^2)_+\bigr)+o_j(1).
 \tag{5.9}
\]

This proves the required second-derivative budget without
differentiating the weight and without an estimate of the anti-self-dual
Weyl curvature.

\hypertarget{a-weak-bianchi-self-test-and-its-incompatible-lower-bound}{%
\subsection{6. A weak Bianchi self-test and its incompatible lower
bound}\label{a-weak-bianchi-self-test-and-its-incompatible-lower-bound}}

\hypertarget{tensor-normalization-and-the-bianchi-upper-bound}{%
\subsubsection{6.1. Tensor normalization and the Bianchi upper
bound}\label{tensor-normalization-and-the-bianchi-upper-bound}}

For an orthonormal basis \(\epsilon_a\) of \(\Lambda^+\), put
\(\omega_a=\sqrt2\epsilon_a\). Associate to a symmetric endomorphism
\(A\) the four-tensor

\[
 W(A)_{ijkl}=\frac12\sum_{a,b}A_{ab}(\omega_a)_{ij}(\omega_b)_{kl}.
 \tag{6.1}
\]

Let

\[
                    J_\psi=\operatorname{div}W(A_\psi),
\]

using the full three-index norm \(\sum_{j,k,l}|J_{jkl}|^2\). Changing
the overall sign of divergence changes none of the following squared
norms.

The contracted Bianchi identity expresses \(\operatorname{div}W^+\) as a
parallel algebraic contraction of \(\nabla\operatorname{Ric}^\circ\) and
\(\nabla s\). This follows by contracting the Riemannian Bianchi
identity, subtracting the Schouten part, and projecting the final
two-form pair onto \(\Lambda^+\); that projection is parallel. Since
\(\nabla\operatorname{Ric}=\nabla E\) and
\(\nabla s=\nabla\operatorname{tr}E\), integration by parts and
\(W(A_\psi)=W^+-W(R)\) show, for every smooth test tensor \(Z\), that

\[
 \left|\int\langle J_\psi,Z\rangle\right|
       \le C(\|E\|_2+\|R\|_2)\|\nabla Z\|_2.
 \tag{6.2}
\]

Set \(Z=J_\psi\). It is smooth, and
\(\|\nabla J_\psi\|_2\le C\|\nabla^2A_\psi\|_2\). Equations (4.6) and
(5.9) therefore give

\[
 \frac1V\int|J_\psi|^2
            \le C\sqrt\kappa\sqrt{1+\kappa+B_a}+o_j(1).
 \tag{6.3}
\]

This is a weak identity tested against a smooth tensor already supplied
by the spinor. It asserts no pointwise bound on
\(\nabla\operatorname{Ric}\).

\hypertarget{the-exact-divergence-identity}{%
\subsubsection{6.2. The exact divergence
identity}\label{the-exact-divergence-identity}}

Define

\[
 B_\psi=\frac{s}{3}\rho^2-2\langle W^+\psi,\psi\rangle,
 \qquad C_\psi=\langle\iota_{\nabla\rho}\psi,*d\psi\rangle.
\]

We now derive the identity used to lower-bound the left side of (6.3).
Pairing the self-dual Hodge--Weitzenböck equation with \(\rho^2\psi\)
and integrating by parts gives

\[
 \int\rho^2\bigl(|\nabla\psi|^2+B_\psi-2|d\psi|^2\bigr)
          =4\int\rho C_\psi-2\int\rho^2|d\rho|^2.
 \tag{6.4}
\]

To see the mixed terms explicitly, use \(d^*\psi=-*d\psi\),
\(d^*(\rho^2\psi)=\rho^2d^*\psi-2\rho\iota_{\nabla\rho}\psi\), and
\(d(\rho^2\psi)=\rho^2d\psi+2\rho d\rho\wedge\psi\). The two
exterior-derivative mixed terms add to \(4\rho C_\psi\); the
rough-Laplacian mixed term is \(2\rho^2|d\rho|^2\).

The required pointwise first-jet identities are

\[
 |\nabla A_\psi|^2
       =\frac{\rho^2}{32}|\nabla\psi|^2
                                +\frac{\rho^2}{96}|d\rho|^2,
 \tag{6.5}
\]

\[
 |\nabla A_\psi|^2-|J_\psi|^2
    =\frac{\rho^2}{64}(|\nabla\psi|^2-2|d\psi|^2)
      +\frac{\rho^2}{32}|d\rho|^2-\frac\rho{16}C_\psi.
 \tag{6.6}
\]

Here is a direct component verification that fixes all normalization
factors. At a point with \(\rho>0\), choose a parallel orthonormal frame
so that

\[
 \begin{aligned}
 \epsilon_1&=(dx^0\wedge dx^1+dx^2\wedge dx^3)/\sqrt2,\\
 \epsilon_2&=(dx^0\wedge dx^2+dx^3\wedge dx^1)/\sqrt2,\\
 \epsilon_3&=(dx^0\wedge dx^3+dx^1\wedge dx^2)/\sqrt2,\\
 \psi&=\rho\epsilon_1,\qquad
 \nabla_m\psi=c_m\epsilon_1+a_m\epsilon_2+b_m\epsilon_3.
 \end{aligned}
\]

Then \(c_m=\partial_m\rho\), and differentiating (4.4) gives

\[
 \nabla_m A_\psi=
 \begin{pmatrix}
  -\rho c_m/6&-\rho a_m/8&-\rho b_m/8\\
  -\rho a_m/8&\rho c_m/12&0\\
  -\rho b_m/8&0&\rho c_m/12
 \end{pmatrix}.                                         \tag{6.7}
\]

The squared matrix norm gives (6.5). For (6.6), insert (6.7) in
\(J_{jkl}=\sum_{m,a,b}(\nabla_mA)_{ab}(\epsilon_a)_{mj}(\epsilon_b)_{kl}\),
use \(\sum_{k,l}(\epsilon_b)_{kl}(\epsilon_c)_{kl}=2\delta_{bc}\), and
compute
\(d\psi_{ijk}=\nabla_i\psi_{jk}+\nabla_j\psi_{ki} +\nabla_k\psi_{ij}\)
in the displayed basis. Explicitly,

\[
 *d\psi=\frac1{\sqrt2}
 \begin{pmatrix}
  -a_2-b_3-c_1\\
   a_3-b_2+c_0\\
   a_0+b_1-c_3\\
  -a_1+b_0+c_2
 \end{pmatrix},
 \qquad
 \iota_{\nabla\rho}\psi
       =\frac{\rho}{\sqrt2}(-c_1,c_0,-c_3,c_2).
\]

These expressions and (6.7) give (6.6) by summing squares and collecting
the mixed terms in \(C_\psi\). Thus (6.5)--(6.6) are algebraic
identities in the twelve arbitrary numbers \(c_m,a_m,b_m\), with no
field equation imposed. At a spinor zero, both \(\psi\) and its first
derivative vanish, and \(d\rho=0\); the identities extend there as well.

Combining (6.4)--(6.6) eliminates \(\int\rho^2|\nabla\psi|^2\) and gives

\[
 \begin{aligned}
 \int|J_\psi|^2={}&-\frac1{64}\int\rho^2 B_\psi
       +\frac1{16}\int\rho^2|d\psi|^2
       -\frac5{96}\int\rho^2|d\rho|^2
       +\frac18\int\rho C_\psi.
 \end{aligned}                                              \tag{6.8}
\]

\hypertarget{the-positive-lower-bound-and-contradiction}{%
\subsubsection{6.3. The positive lower bound and
contradiction}\label{the-positive-lower-bound-and-contradiction}}

Because \(A_\psi\psi=-\rho^2\psi/12\), the curvature term in (6.8) is

\[
 -\frac{\rho^2B_\psi}{64}
        =-\frac{s\rho^4}{192}-\frac{\rho^6}{384}
                                +\frac{\rho^2\langle R\psi,\psi\rangle}{32}.
 \tag{6.9}
\]

At \(s=-4\), \(\rho=\sqrt2\), and \(R=0\), it equals \(1/16\). On the
interval \(0\le\rho\le4\), the error in (6.9) is bounded by

\[
                  C\bigl(|e|+|\rho-\sqrt2|+|R|\bigr).
\]

Also \(|C_\psi|\le\rho|d\rho||d\psi|\). Discard the nonnegative
\(|d\psi|^2\) term in (6.8), apply (3.9), (4.2), and (4.6), and divide
by \(V\). Cauchy--Schwarz gives

\[
 \frac1V\int|J_\psi|^2
                   \ge\frac1{16}-C(\sqrt\kappa+\kappa)+o_j(1).
 \tag{6.10}
\]

For each fixed metric let the smooth-weight approximation errors tend to
zero in (6.3) and (6.10). The two numerical inequalities remain valid
without choosing any limit of the spinors. For a sufficiently small
absolute \(\kappa_0\le1\), the lower bound is at least \(1/32\) when
\(\kappa\le\kappa_0\). The upper bound then forces

\[
                         \kappa(1+\kappa+B_a)\ge c_1>0.
\]

As \(1+\kappa+B_a\le2(1+B_a)\) in this range, (1.8) follows. For
\(\kappa>\kappa_0\), it follows after decreasing the same universal
constant. This proves Theorem 1.3. \(\square\)

\hypertarget{from-the-calibration-gap-to-topological-integral-rigidity}{%
\subsection{7. From the calibration gap to topological integral
rigidity}\label{from-the-calibration-gap-to-topological-integral-rigidity}}

We now prove Theorem 1.2, first at the negative normalization. Assume
(1.5), and put

\[
 r=\frac1V\int|\operatorname{Ric}+g|^2,\qquad
 d=\left(\frac TA-\frac23\right)_+.
\]

Continue to use \(\alpha=\sqrt{2/3}|W^+|\) and
\(w=\lambda_{\min}(W^+)\), and introduce

\[
 t=s+24\alpha=s+8\sqrt6|W^+|,\qquad
 b=2s/3+2w,\qquad
 G=(2s/3-2\alpha)^2-b_-^2.                               \tag{7.1}
\]

Since \(b\ge2s/3-2\alpha\), we have \(G\ge0\), including where the
scalar curvature is positive. Direct expansion gives

\[
 t^2+18G=216\left(s^2/24+2|W^+|^2\right)-18b_-^2.
\]

Apply (1.1) and (2.12) to obtain

\[
                    \int(t^2+18G)\le108P+864\pi^2Ad.
 \tag{7.2}
\]

Here \(p\ge A\) and \(T-2p/3\le T-2A/3\le Ad\). Since
\(\int|E|^2=P+\frac14\int e^2\), (7.2) implies

\[
             \int(|E|^2+e^2+t^2+G)\le109rV+864\pi^2Ad.
 \tag{7.3}
\]

We record the pointwise orbit comparison, including its scalar-positive
branch. Set \(z=w+\alpha\); trace-freeness gives \(0\le z\le\alpha\),
and minimizing over the orbit gives exactly

\[
 F=\frac32(\alpha-1/6)^2+z/2,\qquad
                         \alpha=(t-e+4)/24.              \tag{7.4}
\]

If \(s\le-3\), then \(b<0\) and \(G=z(-8s/3+8\alpha-4z)\ge8z\). Thus

\[
                       F\le(t^2+e^2)/192+G/16.
\]

If \(s>-3\), then \(e>1\). Use \(z\le\alpha\), \(4-e\le3e^2\), and
\(t\le(t^2+1)/2\) in (7.4) to get

\[
                       F\le(t^2+5e^2)/64.
\]

In both cases,

\[
                             F\le e^2+t^2+G.             \tag{7.5}
\]

The second completed square is

\[
 27\left(s^2/24+2|W^+|^2\right)-q^2=t^2/8.
\]

It and (1.1) yield

\[
 \Delta_+\le\frac{27}{2}rV+108\pi^2Ad.
\]

Combining this with (7.3)--(7.5),

\[
                     \kappa\le\frac{245}{2}r
                                      +972\pi^2\frac AVd.
 \tag{7.6}
\]

The scalar estimate (2.13) gives

\[
 32\pi^2A\le\int s^2\le32V+8rV,
              \qquad \frac AV\le\frac{4+r}{4\pi^2}.
 \tag{7.7}
\]

For \(r\le1\), therefore,

\[
                              \kappa\le1215(r+d).        \tag{7.8}
\]

When \(r\le1/4\), reverse triangle inequality in \(L^2\) and (1.1) also
imply

\[
 4\pi^2T
     \ge\left[\frac{(4-2\sqrt r)^2}{24}-\frac r2\right]V
     \ge V/4>0.                                          \tag{7.9}
\]

The square condition \(a^2\ge T\) then gives \(B_a=0\). Apply Theorem
1.3 and (7.8). If \(r>1/4\), no argument is needed. In either case

\[
       r+d\ge\varepsilon_0,
       \qquad \varepsilon_0:=\min\{1/4,c_{\mathrm{cal}}/1215\}.
 \tag{7.10}
\]

For reference, a weaker square hypothesis has a correspondingly
nonuniversal version. If \(a^2\ge-C_AA\), \(C_A\ge0\), the same proof
gives

\[
 r+d\ge\min\left\{1,
          \frac{c_{\mathrm{cal}}}{1215(1+5C_A/(4\pi^2))}\right\}.
 \tag{7.11}
\]

If instead \(a^2\ge T-C_AA\), replace the first entry of that minimum by
\(1/4\), using (7.9). An arbitrary monopole class must not silently be
assumed to have \(a^2\ge T\).

\hypertarget{positive-scale-arbitrary-scale-and-the-ricci-flat-case}{%
\subsection{8. Positive scale, arbitrary scale, and the Ricci-flat
case}\label{positive-scale-arbitrary-scale-and-the-ricci-flat-case}}

\hypertarget{the-positive-normalization}{%
\subsubsection{8.1. The positive
normalization}\label{the-positive-normalization}}

Put \(r_+=V^{-1}\int|\operatorname{Ric}-g|^2\). Since
\(s_-^2\le(s-4)^2\le4|\operatorname{Ric}-g|^2\), (2.13) gives

\[
                            A/V\le r_+/(8\pi^2).
\]

If \(r_+\le1/4\), the analogue of (7.9) gives \(4\pi^2T/V\ge1/4\). If
also \(d\le1/4\), then \(T/A\le2/3+d\), so

\[
  \frac14\le\frac{4\pi^2T}{V}
             \le\frac{r_+}{2}\left(\frac23+d\right)
             \le\frac{11}{96},
\]

a contradiction. In particular \(r_++d>1/4\).

\hypertarget{scale-optimization}{%
\subsubsection{8.2. Scale optimization}\label{scale-optimization}}

For \(\lambda\ne0\), apply the appropriate normalized result to
\(|\lambda|g\). The positive harmonic subspace and the conformal
solvability hypothesis are unchanged. The normalized error becomes

\[
 \frac{\int|\operatorname{Ric}_g-\lambda g|^2}{\lambda^2V}.
\]

Taking the infimum proves (1.6). This completes Theorem 1.2. \(\square\)

For interpretation, write

\[
 R_2=\int|\operatorname{Ric}|^2,\qquad S=\int s,\qquad
 L_s=\int(s-S/V)^2.
\]

When \(R_2>0\), direct minimization in \(1/\lambda\) gives

\[
              \mathcal E(g)=4-\frac{S^2}{VR_2}
                             =\frac{4P+L_s}{R_2}.        \tag{8.1}
\]

This formula includes \(S=0\), when the infimum is \(4\) and need not be
attained. If \(R_2=0\), definition (1.3) gives \(\mathcal E(g)=4\), not
zero. Thus a lower bound for \(\mathcal E\) alone does not exclude a
Ricci-flat metric. Ricci-flat metrics in Theorem 1.2 are excluded
separately by (2.13), since \(p\ge A>0\).

For an Einstein metric with nonzero Einstein constant, \(\mathcal E=0\).
A pointwise operator error also dominates the integral one: for
\(\sigma=\pm1\),

\[
 \delta_\sigma(g):=\|\operatorname{Ric}^{\sharp}-\sigma I\|_{L^\infty,\mathrm{op}}
      \quad\Longrightarrow\quad\mathcal E(g)\le4\delta_\sigma(g)^2.
 \tag{8.2}
\]

There are exactly four Ricci eigenvalues, which explains this factor.

\hypertarget{metricwise-polarized-classes-and-chamber-bookkeeping}{%
\subsection{9. Metricwise polarized classes and chamber
bookkeeping}\label{metricwise-polarized-classes-and-chamber-bookkeeping}}

The analytic theorem needs the quantifiers

\[
 \boxed{\begin{gathered}
   \text{for every }g\text{ there exists a Spin}^c
       \text{ structure }\mathfrak s_g\\
   \text{such that for every }f>0,\text{ equation (1.2)}\\
   \text{has a solution for the same }\mathfrak s_g.
 \end{gathered}}
 \tag{9.1}
\]

It does not require one universal choice of signs for all metrics. This
distinction is essential when \(b_+=1\).

\hypertarget{the-connected-sum-input-with-insertions}{%
\subsubsection{9.1. The connected-sum input with
insertions}\label{the-connected-sum-input-with-insertions}}

Let \(X\) have \(b_+(X)>0\), and suppose a Spin\(^c\) structure
\(\mathfrak s\), with real first Chern class \(c\), has a nonzero
Seiberg--Witten invariant evaluated on a homogeneous insertion \(z\).
The invariant is defined on

\[
 \mathbb A(X)=\mathbb Z[U]\otimes\Lambda^*H_1(X;\mathbb Z)/\mathrm{torsion},
                 \qquad\deg U=2,\quad\deg H_1=1.
\]

Its expected dimension is \((c^2-T_X)/4\), where
\(T_X=2\chi(X)+3\tau(X)\). Assume

\[
                    T_X>0,\qquad A:=c^2\ge T_X.          \tag{9.2}
\]

The second inequality follows automatically from the nonnegative degree
of a nonzero ordinary invariant, but is displayed because it is used in
the proof. For \(b_+=1\), the nonzero invariant here must be in the
small-perturbation chamber at the relevant polarization, not merely in
some unspecified chamber.

Let

\[
 M=X\#k\,\overline{\mathbb{CP}}^{2}\#\ell(S^1\times S^3),
 \qquad D=k+4\ell,\qquad T_M=T_X-D.                      \tag{9.3}
\]

For a prescribed positive polarization \(H=H_g^+\), choose the signs of
the exceptional generators \(E_j\) so that

\[
                         c_H^+\mathbin\cdot E_j\le0,
\]

and set

\[
                             a_H=c-\sum_{j=1}^kE_j.
 \tag{9.4}
\]

Then

\[
 \begin{aligned}
 (a_H^+)^2
   &=(c_H^+)^2-2\sum_jc_H^+\mathbin\cdot E_j
                              +\left|\sum_j(E_j)_H^+\right|^2
      \ge(c_H^+)^2\ge c^2=A,\\
 a_H^2&=A-k\ge T_X-k-4\ell=T_M.
 \end{aligned}                                              \tag{9.5}
\]

The ordinary blow-up formula with exceptional coefficient \(\pm1\)
preserves expected dimension and preserves the original insertion \(z\).
The \(S^1\times S^3\) operation raises expected dimension by one and
adds the degree-one insertion of its new circle. After \(\ell\)
operations the nonzero invariant is therefore evaluated on

\[
                           z\,\mu(C_1)\cdots\mu(C_\ell).
 \tag{9.6}
\]

These statements include positive-dimensional base invariants. They are
the special cases of the blow-up and surgery formulas in
\href{https://arxiv.org/pdf/math/0005268}{Ozsváth--Szabó, Theorem 3.2,
Proposition 2.2 and Remark 2.3}. The last remark supplies the
chamber-compatible version when \(b_+=1\). Using only a zero-dimensional
numerical invariant and omitting the insertions in (9.6) would not
justify the general statement.

\hypertarget{what-the-chamber-condition-means}{%
\subsubsection{9.2. What the chamber condition
means}\label{what-the-chamber-condition-means}}

For \(b_+(X)=1\), choose a unit forward generator \(\omega_H\) of the
positive line on \(M\) so that \(c\cdot\omega_H>0\). The \(X\) component
\(\omega_X\) has positive square, lies in the corresponding forward
cone, and satisfies \(c\cdot\omega_X>0\), because \(c^2>0\). The chosen
exceptional signs give \(E_j\cdot\omega_H\le0\), hence

\[
                          a_H\cdot\omega_H>0.             \tag{9.7}
\]

Thus the small-perturbation discriminants for the base and the selected
blow-up structure have the same sign in corresponding chamber
conventions. Adding \(S^1\times S^3\) leaves the second-cohomology
polarization unchanged. The gluing formulas consequently give a nonzero
invariant in the chamber containing zero perturbation for this selected
class, provided that was the specified nonzero base chamber.

For any metric with this polarization, take generic perturbations
tending to zero in that chamber. Nonzero invariant gives solutions, and
the usual fixed-metric Seiberg--Witten compactness gives an unperturbed
solution. A reducible limit would have \(F_A^+=0\), contradicting
\((a_H^+)^2\ge A>0\). Every metric conformal to \(g\) has the same
polarization and the same nonzero discriminant sign. Thus the
\emph{same} class (9.4) satisfies Definition 1.1 for every weight, as
required by (9.1).

For \(b_+>1\), there is no chamber ambiguity in this argument. For
\(b_+=1\), the statement must not be replaced by the assertion that
every exceptional sign choice is an ordinary all-metric monopole class.

\hypertarget{the-minimal-general-type-case}{%
\subsubsection{9.3. The minimal general-type
case}\label{the-minimal-general-type-case}}

If \(X\) is a minimal compact complex surface of general type with its
complex orientation, then

\[
                          A=c_1(X)^2=T_X>0.
\]

The canonical Spin\(^c\) structure has a nonzero Seiberg--Witten
invariant. For \(b_+>1\), the standard Kähler/symplectic nonvanishing
statement supplies the input directly. For \(b_+=1\), the canonical
class \(K\) has \(K^2>0\) and positive pairing with a Kähler form. It is
therefore in the forward time-like cone and pairs positively with every
forward period ray. For \(c=-K\), the small-perturbation chamber agrees
with the canonical nonzero chamber, up to the simultaneous choice of
sign convention for the chamber discriminant. This is the
light-cone/chamber argument underlying
\href{https://arxiv.org/pdf/math/0003068}{LeBrun's polarized
construction, Proposition 3.2}. The chamber convention and canonical
nonvanishing can also be read in
\href{https://arxiv.org/pdf/math/0012048}{Li--Liu, §2.1 and Theorem
2.2}.

The positivity of \(K\) on the symplectic or Kähler form matters:
\(b_+=1\) and \(K^2>0\) alone do not identify the nonzero
small-perturbation chamber for an arbitrary symplectic manifold.

\hypertarget{the-einstein-obstruction-and-its-precise-relation-to-lebrun}{%
\subsection{10. The Einstein obstruction and its precise relation to
LeBrun}\label{the-einstein-obstruction-and-its-precise-relation-to-lebrun}}

\textbf{Theorem 10.1 (connected-sum integral obstruction).} Under the
gauge-theoretic and chamber assumptions of §9.1, every metric on (9.3)
satisfies

\[
 \mathcal E(g)+\left(\frac{T_M}{A}-\frac23\right)_+
                         \ge\varepsilon_0.               \tag{10.1}
\]

It also satisfies the slightly weaker but base-topological statement

\[
 \boxed{\displaystyle
 \mathcal E(g)+\left(\frac13-\frac{D}{T_X}\right)_+
                         \ge\varepsilon_0.}              \tag{10.2}
\]

For minimal general type, \(A=T_X=c_1(X)^2\), so the two statements
coincide.

\textbf{Proof.} Equation (9.5) and the conformal solvability just proved
give all the hypotheses of Theorem 1.2. If \(T_M>0\), then \(A\ge T_X\)
implies \(T_M/A\le T_M/T_X\), yielding (10.2). If \(T_M\le0\), both
positive-part terms are zero. This case split avoids reversing an
inequality by dividing a negative numerator by the larger denominator.
\(\square\)

In particular, put \(\gamma=\varepsilon_0/2\). Then

\[
             D\ge(1/3-\gamma)T_X
 \quad\Longrightarrow\quad
             \mathcal E(g)\ge\varepsilon_0/2
             \quad\hbox{for every metric }g.             \tag{10.3}
\]

There is no Einstein metric in this band: nonzero Einstein constants
contradict (10.3), and zero Einstein constant contradicts (2.13). At
either normalized sign, (8.2) also gives

\[
                    \delta_\sigma(g)\ge\sqrt{\varepsilon_0/8}.
 \tag{10.4}
\]

An Einstein metric would in fact require \(D/T_X\le1/3-\varepsilon_0\).
The equality endpoint of that necessary condition is not excluded by
(10.2) alone; the closed band (10.3) avoids that endpoint issue.

\hypertarget{exact-source-scope}{%
\subsubsection{10.2. Exact source scope}\label{exact-source-scope}}

The version of \href{https://arxiv.org/pdf/math/0003068}{LeBrun's paper
verified here} is arXiv:math/0003068v1, March 2000. Its Proposition 3.2
constructs the metricwise polarized classes on the connected sums in
(9.3). Its Theorem 3.3 states the classical threshold

\[
                              k+4\ell\ge T_X/3,
\]

and Corollary 3.4 specializes the result to blow-ups of minimal complex
surfaces of general type. We use the ordinary nontrivial Seiberg--Witten
input, including its insertion and chamber meaning, not merely the
topology of a manifold with the same Euler characteristic and signature.
The proof above handles the positive-budget case of the general theorem,
not only its general-type corollary.

The numerical band below the classical threshold is nonvacuous in
families with unbounded budget. For example, let
\(X_m=\Sigma_2\times\Sigma_m\), \(m\ge2\). Then \(A=T_{X_m}=8(m-1)\).
With \(\ell=0\) and \(k=\lceil A/3\rceil-1\),

\[
                       0<1/3-k/A\le1/A\longrightarrow0.
\]

For sufficiently large \(m\), these ratios lie in (10.3). Since the
universal analytic constant has not been made numerical, this argument
does not specify the first qualifying \(m\).

\hypertarget{nonpositive-budgets-and-the-exceptional-zero-corner}{%
\subsubsection{10.3. Nonpositive budgets and the exceptional zero
corner}\label{nonpositive-budgets-and-the-exceptional-zero-corner}}

There is a direct all-metric formula, independent of monopoles. For
\(R_2>0\), (1.1) and (8.1) give

\[
 \mathcal E(g)
   =1+\frac{12\int|W^+|^2-24\pi^2T+L_s}{R_2}.           \tag{10.5}
\]

Thus \(T\le0\) implies \(\mathcal E(g)\ge1\). For a Ricci-flat metric
its value is \(4\) by definition, so the same lower bound holds. This
lower bound is not a Ricci-flat exclusion.

For an Einstein metric, (1.1) reduces to

\[
               4\pi^2T=\int(s^2/24+2|W^+|^2)\ge0.
\]

Consequently, when \(T_X<0\), or when \(T_X=0\) and \(D>0\), the
connected sum has \(T_M<0\) and admits no Einstein metric. These are
classical Chern--Gauss--Bonnet consequences and require no
positive-budget monopole argument.

When \(T_X=D=0\), however, equality permits Ricci-flat, anti-self-dual
metrics in the chosen orientation. K3 and the four-torus provide the
standard zero-budget examples. The unqualified statement of Theorem 3.3
in the verified preprint includes this corner, but its appeal to the
Hitchin--Thorpe inequality does not exclude these equality cases. We
therefore do \textbf{not} claim an unqualified subsumption of that
literal zero-corner wording, nor infer an exception-free statement for a
separately unverified published version. The conclusions established
here are: the positive-budget integral strengthening with the stated
gauge input; the classical negative-\(T_M\) exclusion; and an explicit
reservation of the Ricci-flat zero corner.

\hypertarget{logical-dependencies-and-limitations}{%
\subsection{11. Logical dependencies and
limitations}\label{logical-dependencies-and-limitations}}

The analytic proof has the following complete chain:

\[
 \begin{gathered}
 \text{one actual conformal Spin}^c\text{ structure}
 \ \Longrightarrow\ \text{solutions for smooth truncated weights}\\
 \Longrightarrow\ \rho\le4\text{ and small retained first-order defects}
 \ \Longrightarrow\ W^+\approx A_\psi\text{ with }A_\psi\text{ smooth}\\
 \Longrightarrow\ \text{bundle Hessian budget and weak Bianchi upper bound}\\
 \Longrightarrow\ \text{incompatibility with the exact positive divergence term}.
 \end{gathered}
\]

The proof neither differentiates an eigenprojection nor normalizes
\(\psi\) across its zero set. Smooth approximation of weights is used
only for numerical inequalities at a fixed metric; there is no claimed
uniform regularity or limiting solution for a nonsmooth weight. The
Hessian estimate uses the positive-spinor connection, which is why no
\(L^2\) budget for \(W^-\) enters. The topological theorem additionally
needs the class-square bound and the actual class/chamber construction;
the scalar curvature inequality by itself is not a substitute for those
inputs.

This argument proves an obstruction and an integral gap. It does not
prove convergence of near-Einstein metrics, classify all equality
metrics of other curvature functionals, produce an Einstein metric, or
settle arbitrary pointwise Ricci-pinching problems without the
monopole/topological hypotheses displayed above.

\hypertarget{references-and-provenance}{%
\subsection{References and provenance}\label{references-and-provenance}}

\begin{itemize}
\tightlist
\item
  Claude LeBrun, \emph{Ricci Curvature, Minimal Volumes, and
  Seiberg--Witten Theory},
  \href{https://arxiv.org/pdf/math/0003068}{arXiv:math/0003068v1},
  especially Proposition 3.2, Theorem 3.3 and Corollary 3.4. The source
  version and the zero-budget caveat are specified in §10.
\item
  Peter Ozsváth and Zoltán Szabó, \emph{Higher Type Adjunction
  Inequalities in Seiberg--Witten Theory},
  \href{https://arxiv.org/pdf/math/0005268}{arXiv:math/0005268},
  especially Proposition 2.2, Remark 2.3 and Theorem 3.2. These supply
  insertion-aware surgery and blow-up formulas.
\item
  Tian-Jun Li and Ai-Ko Liu, \emph{Uniqueness of Symplectic Canonical
  Class, Surface Cone and Symplectic Cone of 4-Manifolds with
  \(b^+=1\)},
  \href{https://arxiv.org/pdf/math/0012048}{arXiv:math/0012048}, §2.1
  and Theorem 2.2, for chamber conventions and canonical nonvanishing.
\item
  Claude LeBrun, \emph{Four-Manifolds, Curvature Bounds, and Convex
  Geometry}, \href{https://arxiv.org/pdf/math/0611450}{arXiv preprint},
  especially Lemma 3.7, for the conformal weighted-equation input.
\end{itemize}

This text consolidates the local weighted-monopole,
integral-calibration, all-sign, chamber-selection and LeBrun-comparison
arguments. Those earlier working records are provenance, not proof
dependencies of the analytic argument presented here. Numerical and
symbolic checks are useful for detecting transcription errors; none is
used in place of a mathematical identity or a gauge-theoretic existence
theorem.

END COMPLETE SOURCE

\begin{center}\rule{0.5\linewidth}{0.5pt}\end{center}

\hypertarget{source-2-round21_analytic.md}{%
\subsection{Source 2:
round21\_analytic.md}\label{source-2-round21_analytic.md}}

Original path: mathematics/current/round21\_analytic.md

SHA-256:
9e39e932ae9f109d9ab3f959a243046e243beb20900c3f4e95be8dc424f5eb22

BEGIN COMPLETE SOURCE

\hypertarget{integral-rigidity-for-bounded-self-dual-forms}{%
\section{Integral rigidity for bounded self-dual
forms}\label{integral-rigidity-for-bounded-self-dual-forms}}

This note isolates the analytic mechanism behind the integral monopole
gap. It proves an abstract theorem whose input is a bounded self-dual
two-form with a Hodge-Laplacian budget, rather than a monopole spinor.
This makes the production of a calibration and its differential budget
separate inputs. It is an analytic criterion; no new family of smooth
topologies is asserted to satisfy its hypotheses.

All manifolds are smooth, closed, oriented and four-dimensional.
Curvature is normalized relative to \(\operatorname{Ric}=-g\). Two-forms
have the exterior-algebra norm; \(W^+\) has the Hilbert--Schmidt
operator norm. Write \(\fint=\operatorname{Vol}^{-1}\int\), and let
\(\Delta_H\) denote the nonnegative Hodge Laplacian.

\hypertarget{the-bounded-two-form-criterion}{%
\subsection{1. The bounded two-form
criterion}\label{the-bounded-two-form-criterion}}

For a smooth self-dual two-form \(\psi\), set

\[
 q=|\psi|^2,\qquad
 A_\psi={q\over24}I-{1\over8}\psi\otimes\psi,
 \qquad R=W^+-A_\psi.
\]

Define the defect and source budget

\[
 d=\fint_M\left(|\operatorname{Ric}+g|^2+|R|^2+(q-2)^2\right),
 \qquad B=\fint_M|\Delta_H\psi|^2.
\]

\textbf{Theorem.} For each finite \(K>0\), there is \(c(K)>0\) such that
\(\|\psi\|_\infty\le K\) implies

\[
                    \boxed{(1+B)d\ge c(K).}              \tag{1}
\]

\textbf{Small-source corollary.} For a possibly smaller \(c_0(K)>0\),

\[
 \boxed{\quad
 \fint_M\left(
 |\operatorname{Ric}+g|^2+|R|^2+(q-2)^2+|\Delta_H\psi|^2
 \right)\ge c_0(K).
 \quad}                                                     \tag{1a}
\]

No spin structure, monopole class, curvature bound, scalar sign, Sobolev
constant, injectivity radius or diameter bound is required. In
particular, for bounded \emph{harmonic} self-dual forms \(B=0\), giving
a gap for the three zeroth-order defects alone. No small
first-derivative defect is assumed, even in the nonharmonic case. The
boundedness of the form and its Hodge-Laplacian control are genuine
hypotheses, not consequences of a cohomology class alone.

The theorem excludes a specific incompatible calibration: Ricci
curvature near \(-g\), Weyl curvature near \(A_\psi\), and \(|\psi|^2\)
near two. For a product of curvature-minus-one surfaces the simple Weyl
eigenvalue is \(-2/3\), not the target \(-1/6\). Thus the theorem does
not exclude the realizable hyperbolic-product Weyl model and does not
assert a universal Ricci gap without the other summands.

\hypertarget{re-derivation-of-the-bundle-hessian-estimate}{%
\subsection{2. Re-derivation of the bundle Hessian
estimate}\label{re-derivation-of-the-bundle-hessian-estimate}}

Let \(u\) be a smooth section of a metric bundle, with compatible
connection \(\nabla\). Set \(H_{ij}=\nabla_i\nabla_j u\), including the
cotangent Levi-Civita correction, \(L=\operatorname{tr}H\), and
\(F_{ij}=[\nabla_i,\nabla_j]\). Use the component convention
\(|F|^2=\sum_{ij}|F_{ij}|_{\rm HS}^2\).

First, integration of the divergence of
\(|\nabla u|^2\langle u,\nabla u\rangle^\sharp\) gives

\[
 \|\nabla u\|_4^2\le(2+\sqrt4)\|u\|_\infty\|H\|_2.
                                                               \tag{2}
\]

This is a covariant integration-by-parts inequality, not a Sobolev
embedding. To check the curvature terms, first interchange the two
Hessian slots:

\[
 \int|H|^2-\sum_{ij}\int\langle H_{ij},H_{ji}\rangle
 =\frac12\sum_{ij}\int|F_{ij}u|^2.
\]

Integrate the second expression once and then commute derivatives on the
bundle-valued one-form \(\nabla u\). The cotangent curvature gives
Ricci; the bundle curvature acts on \(\nabla u\). Integrating the trace
term once more yields

\[
\begin{aligned}
 \int|H|^2={}&\int|L|^2
 -\int\operatorname{Ric}_{ij}\langle\nabla_i u,\nabla_j u\rangle\\
 &+\sum_{ij}\int\langle F_{ij}\nabla_j u,\nabla_i u\rangle
 +\frac12\sum_{ij}\int|F_{ij}u|^2.                         
\end{aligned}\tag{3}
\]

There is no differentiated-curvature term. Swapping the Hessian slots
\emph{before} commuting is important. The scalar specialization fixes
the Ricci sign. With \(K_u=\|u\|_\infty\), Holder, (2), and Young give

\[
 \|\nabla^2u\|_2^2\le
 2\|L\|_2^2+32K_u^2\|\operatorname{Ric}\|_2^2
                         +33K_u^2\|F\|_2^2.              \tag{4}
\]

This recovers the inherited estimate, including its coefficient and the
absence of a hidden Sobolev or curvature-derivative hypothesis.

\hypertarget{proof-of-the-two-form-criterion}{%
\subsection{3. Proof of the two-form
criterion}\label{proof-of-the-two-form-criterion}}

Assume \(d\le1\); otherwise the conclusion is immediate after reducing
\(c(K)\). Put \(\mathcal B=1+B+d\). All constants below depend only on
\(K\) and the dimension. Put \(E=\operatorname{Ric}+g\) and
\(\eta=\Delta_H\psi\).

\hypertarget{the-self-dual-connection-supplies-the-second-order-budget}{%
\subsubsection{3.1 The self-dual connection supplies the second-order
budget}\label{the-self-dual-connection-supplies-the-second-order-budget}}

The curvature of \(\Lambda^+\) is the representation of the Riemann
operator projected to its self-dual output. Hence

\[
 |F^{\Lambda^+}|^2\le C\left(|W^+|^2+s^2+
                               |\operatorname{Ric}^{\circ}|^2\right).
                                                               \tag{5}
\]

The \(W^-\) block does not act on \(\Lambda^+\). Since
\(|A_\psi|^2=q^2/96\) and \(q\le K^2\), the \(E,R\) budgets give

\[
 \fint\left(|\operatorname{Ric}|^2+|F^{\Lambda^+}|^2\right)
 \le C_K(1+d).
\]

The self-dual Weitzenbock equation is

\[
 \nabla^*\nabla\psi=\eta+2W^+\psi-(s/3)\psi.
\]

Apply (2)--(4) to \(u=\psi\). Differentiating the quadratic polynomial
\(A_\psi\) and the scalar \(q\) then gives

\[
 \fint\left(|\nabla^2\psi|^2+|\nabla\psi|^4+
              |\nabla^2A_\psi|^2+|\Delta q|^2\right)
 \le C_K\mathcal B.
                                                               \tag{6}
\]

Here \(\Delta=d^*d\) on functions, and
\(\Delta q=2\langle\nabla^*\nabla\psi,\psi\rangle -2|\nabla\psi|^2\).
All estimates are extensive before division by volume; no lower volume
bound enters.

\hypertarget{amplitude-control-and-a-coercive-divergence-identity}{%
\subsubsection{3.2 Amplitude control and a coercive divergence
identity}\label{amplitude-control-and-a-coercive-divergence-identity}}

The amplitude itself need not have a smooth square root at its zeros.
Use the globally smooth quantity \(q\), rather than divide by
\(|\psi|\). Integration by parts and (6) give

\[
 \fint|dq|^2=\fint(q-2)\Delta q
                  \le C_K\sqrt{d\mathcal B}.             \tag{7}
\]

There is a useful pointwise identity, valid for arbitrary smooth
self-dual forms without first-order equations. Let
\(J=\operatorname{div}W(A_\psi)\). The full contraction in the checker
gives

\[
 |J|^2={q|\nabla\psi|^2\over64}+{q|d\psi|^2\over32}
       -{|dq|^2\over192}
       +{\langle\iota_{\nabla q}\psi,*d\psi\rangle\over32}.
                                                               \tag{8}
\]

The last inner product is at most \(\sqrt q\,|dq|\,|d\psi|\) in absolute
value. Young's inequality therefore gives the pointwise coercivity

\[
 \boxed{\quad
 q\bigl(|\nabla\psi|^2+|d\psi|^2\bigr)
        \le64|J|^2+{4\over3}|dq|^2.
 \quad}                                                     \tag{8a}
\]

This is polynomial at zeros, and introduces no division by \(q\). It
will control the derivative terms for the two-sided selection theorem in
Section 6. Small amplitude error and a second-order budget now supply
the weighted derivative information previously obtained from monopoles.

\hypertarget{weak-bianchi-gives-a-small-self-test}{%
\subsubsection{3.3 Weak Bianchi gives a small
self-test}\label{weak-bianchi-gives-a-small-self-test}}

Let \(W(A)\) be the algebraic Weyl four-tensor represented by \(A\), and
\(J=\operatorname{div}W(A_\psi)\). The projections to \(\Lambda^+\) are
parallel. The contracted Bianchi identity, followed by integration by
parts, gives for every smooth test tensor \(T\)

\[
 \left|\int\langle J,T\rangle\right|
 \le C\left(\|E\|_2+\|R\|_2\right)\|\nabla T\|_2.      \tag{9}
\]

This is a weak derivative estimate, not an estimate on \(\nabla E\). The
tensor \(J\) is smooth, so \(T=J\) is an admissible test. Moreover
\(|\nabla J|\le C|\nabla^2A_\psi|\). Therefore (6) implies

\[
                 \fint|J|^2\le C_K\sqrt{d\mathcal B}.    \tag{10}
\]

\hypertarget{the-same-self-test-has-a-positive-lower-bound}{%
\subsubsection{3.4 The same self-test has a positive lower
bound}\label{the-same-self-test-has-a-positive-lower-bound}}

The inherited weighted Hodge and divergence identities admit the
following entirely polynomial version:

\[
\begin{aligned}
 \int|J|^2=\int\Big(&-{s q^2\over192}-{q^3\over384}
       +{q\langle R\psi,\psi\rangle\over32}
       +{q|d\psi|^2\over16}-{5|dq|^2\over384}\\
       &+{\langle\iota_{\nabla q}\psi,*d\psi\rangle\over16}
                           \Big).                         
\end{aligned}\tag{11}
\]

It follows by testing the self-dual Hodge equation against \(q\psi\) and
contracting the derivative of \(A_\psi\); the pointwise contraction is
independently checked in the accompanying program. Because (11) is
polynomial in \(\psi,\nabla\psi\), it extends across all zeros of
\(\psi\); no spinor-square representation is used.

At \(s=-4,q=2,R=0\), its algebraic part equals \(1/16\). On
\(0\le q\le K^2\), polynomial Lipschitz estimates give an averaged error
bounded by \(C_K\fint(|s+4|+|q-2|+|R|)\le C_K\sqrt d\). The last term is
bounded below by \(-\{q|d\psi|^2+|dq|^2\}/32\). Absorb its first
contribution into the positive \(q|d\psi|^2/16\) already present in
(11). The remaining derivative error is at worst \(-17|dq|^2/384\),
controlled by (7). In particular this argument does not assume \(d\psi\)
small. Consequently

\[
           \fint|J|^2\ge {1\over16}-C_K\sqrt{d\mathcal B}.\tag{12}
\]

Combining (10) and (12) yields \(1/16\le C'_K\sqrt{d(1+B+d)}\). For
\(d\le1\), \(1+B+d\le2(1+B)\), proving (1) after choosing a sufficiently
small \(c(K)>0\). For (1a), if its left side is \(z\le1\), then
\(d\le z\), \(B\le z\), and \((1+B)d\le2z\). Thus
\(c_0(K)=\min\{1,c(K)/2\}\) suffices.

\hypertarget{what-the-monopole-proof-actually-requires}{%
\subsection{4. What the monopole proof actually
requires}\label{what-the-monopole-proof-actually-requires}}

For the inherited integral gap, one fixes a metric \(g\) and a single
spin-c structure with \(p=(a_g^+)^2>0\). The analytic construction only
requires ordinary unperturbed Seiberg--Witten solutions for every metric
in \([g]\). Indeed, solving on \(\widehat g=f^{-2}g\) gives on \(g\)

\[
 D_A\Phi=0,\qquad -iF_A^+=f\sigma(\Phi),
\]

for every positive smooth weight \(f\). The self-dual polarization is
conformally invariant. Every remaining Hessian, Hodge and Bianchi step
takes place on this one metric. Neither \(b_+>1\) nor persistence over
metrics outside \([g]\) is used in those analytic steps.

This is already the scope of \href{round13_chamber.md}{the Round13
chamber result}, not a new Round21 extension. In \(b_+=1\) applications,
chamber invariance and passage from perturbed invariants to unperturbed
solutions still have to be established. The Hessian identity cannot
supply them.

For the monopole argument specifically, the determinant connection
energy is controlled by

\[
 \|F_A\|_2^2=2\|F_A^+\|_2^2-4\pi^2a^2.
\]

That is the origin of the negative-square cost. The theorem in Section 1
avoids this connection by acting directly on \(\Lambda^+\), but replaces
the monopole PDE with a bounded-form and Hodge-budget hypothesis. It
therefore does \textbf{not} remove the square cost from the existing
monopole theorem: the harmonic projection need not preserve a maximum
bound, and a small \(d\psi\) does not bound \(\Delta_H\psi\) by itself.

More precisely, let \(\kappa\) and \(B_a=(-a^2)_+/V\) have the meanings
in \href{round4_rigidity.md}{Round4, Theorem 1}. The inherited
bounded-spinor, Hessian and fourth-power-gradient estimates imply

\[
 B=\fint|\Delta_H\psi|^2\le C(1+\kappa+B_a)+o(1).
\]

Indeed, the quadratic map from \(\Phi\) to \(\psi\) gives
\(|\nabla^2\psi|\le C(|\Phi||\nabla_A^2\Phi|+|\nabla_A\Phi|^2)\); the
two-form Weitzenbock formula adds only the already controlled \(W^+,s\)
terms. The Weyl approximation and amplitude estimates give
simultaneously

\[
                         d\le C\kappa+o(1).
\]

The bounded-form theorem therefore returns
\(\kappa(1+\kappa+B_a)\ge c\), and hence the inherited
\((1+B_a)\kappa\ge c'\) after splitting \(\kappa\le1\) and \(\kappa>1\).
This is a structural recovery of the earlier calibration gap, not a
second independent topological theorem. It displays precisely where a
different construction of a bounded form could replace the monopole
input, and which quantitative budget that construction must supply.

\hypertarget{research-use-and-verification-scope}{%
\subsection{5. Research use and verification
scope}\label{research-use-and-verification-scope}}

The new analytic formulation offers a distinct input route: a family of
bounded, almost-harmonic self-dual forms with the indicated calibration
would be impossible even without monopole solvability. For a harmonic
form, only its maximum bound and the three zeroth-order calibration
defects remain. Producing these hypotheses on a new topological family
is an open step; a formal cohomology vector does not suffice.

The checker verifies the Hessian curvature coefficients on a
noncommuting constant connection, the full twelve-variable divergence
contraction, and the polynomial conversion used in (11). The general
integration-by-parts and PDE arguments are the proofs above, not
consequences of running the checker. The inherited statements have been
checked at these analytic interfaces; this note does not certify every
earlier topological input.

\hypertarget{quantitative-selection-of-the-kuxe4hler-calibration}{%
\subsection{6. Quantitative selection of the Kähler
calibration}\label{quantitative-selection-of-the-kuxe4hler-calibration}}

The forbidden Weyl model is one member of a family, not an isolated
numerical accident. For a constant real parameter \(t\), put

\[
 A_{\psi,t}=t\left({q\over3}I-\psi\otimes\psi\right),
 \quad R_t=W^+-A_{\psi,t},\quad
 d_t=\fint\left(|E|^2+|R_t|^2+(q-2)^2\right).
                                                               \tag{13}
\]

Thus \(t=1/8\) is the model in Theorem (1). At \(q=2\), the Weyl
eigenvalues prescribed by \(A_{\psi,t}\) are \((-4t/3,2t/3,2t/3)\). The
negative Kähler--Einstein value is \(t=1/2\).

\textbf{Branch-selection theorem.} Fix \(K,T<\infty\). For
\(\|\psi\|_\infty\le K\) and \(|t|\le T\),

\[
 \boxed{\quad |t|\,|1-2t|\le C(K,T)
                  \sqrt{d_t(1+B+d_t)}.\quad}             \tag{14}
\]

A quantitative spectral-selection consequence is

\[
 \boxed{\quad
 \operatorname{dist}(t,\{0,1/2\})^2
       \le C(K,T)(1+B)d_t\qquad(d_t\le1).
 \quad}                                                   \tag{14a}
\]

Indeed, \(\operatorname{dist}(t,\{0,1/2\})\le2|t|\,|1-2t|\) for every
real \(t\), since the larger of the distances to zero and one-half is at
least one-quarter. Also \(1+B+d_t\le2(1+B)\). The two branches are
vanishing Weyl amplitude and the negative Kähler--Einstein amplitude,
not an arbitrary list of forbidden constants.

Consequently \((1+B)d_t\) has a uniform positive lower bound when \(t\)
is restricted to a compact subset of \(\mathbb R\setminus\{0,1/2\}\). On
any parameter set \(t_0\le|t|\le T\), with \(t_0>0\), one further has

\[
 |2t-1|+\fint\left(|\nabla\psi|^2+|d\psi|^2\right)
 \le C(K,t_0,T)\sqrt{d_t(1+B+d_t)}\qquad(d_t\le1).       \tag{15}
\]

Thus a bounded-source-budget calibration with vanishing defect and
nonzero Weyl amplitude selects the Kähler value and an almost-parallel
form. This is an integral conclusion; it does not supply coordinates,
pointwise metric closeness or compactness.

\textbf{Proof.} Let \(J_t=\operatorname{div}W(A_{\psi,t})\). Since
\(A_{\psi,t}=8tA_\psi\), one has \(J_t=8tJ\), with \(J\) here denoting
the divergence of the original polynomial on the same form and metric.
The second-order estimates remain valid uniformly for bounded \(t\),
with \(\mathcal B_t=1+B+d_t\). In the weak Bianchi test the derivative
of \(J_t\) has one factor \(|t|\), so

\[
                  \fint|J_t|^2\le C_{K,T}\,|t|
                                      \sqrt{d_t\mathcal B_t}. \tag{16}
\]

The weighted Hodge contraction becomes exactly

\[
\begin{aligned}
 \int|J_t|^2=\int\Big(&-{t^2s q^2\over3}
       -{4t^3q^3\over3}+2t^2q\langle R_t\psi,\psi\rangle
       +4t^2q|d\psi|^2-{5t^2\over6}|dq|^2\\
       &+4t^2\langle\iota_{\nabla q}\psi,*d\psi\rangle
                    \Big).                               
\end{aligned}\tag{17}
\]

Its algebraic part at \(s=-4,q=2,R_t=0\) is \(P(t)=(16/3)t^2(1-2t)\).
The pointwise coercivity (8a) rescales to

\[
 t^2q\bigl(|\nabla\psi|^2+|d\psi|^2\bigr)
          \le |J_t|^2+{4\over3}t^2|dq|^2.                \tag{18}
\]

In (17), the absolute derivative contribution is at most
\(6|J_t|^2+(65/6)t^2|dq|^2\): the positive closedness term costs at most
\(4|J_t|^2+(16/3)t^2|dq|^2\), and Young's inequality bounds the cross
term by \(2|J_t|^2+(14/3)t^2|dq|^2\). This uses (18), not an assumption
that \(d\psi\) is small. Together with the zeroth-order algebraic error
and the amplitude estimate (7),

\[
 |P(t)|\le 7\fint|J_t|^2+
                  C_{K,T}t^2\sqrt{d_t\mathcal B_t}.       \tag{18a}
\]

This is valid on both sides of \(t=1/2\), including negative \(t\).
Combine (16) and (18a), divide by \(|t|>0\), and use \(|t|\le T\),
obtaining (14). The case \(t=0\) is automatic. For \(d_t>1\), (14)
follows after enlarging the constant since its left side is bounded and
its right square root is at least one.

For (15), (14) first bounds \(|2t-1|\) by its right side divided by
\(t_0\). On \(\{q\ge1\}\), (18), (16) and (7) give an averaged bound
\(C_{K,t_0,T}\sqrt{d_t\mathcal B_t}\) for \(|\nabla\psi|^2+|d\psi|^2\).
On \(\{q<1\}\), the relative volume is at most \(\fint(q-2)^2\le d_t\).
The fourth-power gradient estimate (6), and
\(|d\psi|^2\le3|\nabla\psi|^2\), bound its contribution by
\(C_{K,T}\sqrt{d_t\mathcal B_t}\). This proves (15) without assuming
closedness.

For comparison, the unweighted Hodge identity is

\[
 \fint|\nabla\psi|^2
 =2\fint|d\psi|^2-
  \fint\left({s q\over3}+{4tq^2\over3}
                               -2\langle R_t\psi,\psi\rangle\right).
                                                               \tag{19}
\]

The parenthesized expression differs in average from \((8/3)(2t-1)\) by
at most \(C_{K,T}\sqrt{d_t}\). Nonnegativity of the left side gives the
one-sided estimate

\[
 (2t-1)_+\le C_{K,T}\left(\sqrt{d_t}+\fint|d\psi|^2\right).
                                                               \tag{20}
\]

It becomes source-budget-free when \(d\psi=0\), but not for arbitrary
forms under zeroth-order control alone.

At zero defect and \(t\ne0\), (14) forces \(t=1/2\), and (15) gives
\(\nabla\psi=0\). The norm-squared-two self-dual form is then the
parallel Kähler form of an orthogonal complex structure, and
\(\operatorname{Ric}=-g\) makes the metric negative Kähler--Einstein.
This endpoint is genuinely realized; parallelness of the Kähler form
does not imply that the metric is a global product. At \(t=0\), the
divergence polynomial loses sensitivity; (14) makes no assertion that an
exact calibrated example exists there.

This extension identifies the realizable equality model of the analytic
mechanism. It does not change which monopole classes are available, and
it does not by itself extend the topological range of the earlier
integral gap.

\hypertarget{the-same-analytic-criterion-for-either-einstein-sign}{%
\subsection{7. The same analytic criterion for either Einstein
sign}\label{the-same-analytic-criterion-for-either-einstein-sign}}

The spectral-selection argument itself is not restricted to the negative
Einstein sign. For \(\sigma\in\{-1,+1\}\), put

\[
E_\sigma=\operatorname{Ric}-\sigma g,\qquad
d_{\sigma,t}=\fint\left(
|E_\sigma|^2+|W^+-A_{\psi,t}|^2+(q-2)^2\right).
\]

With the same bounded-form and bounded-parameter hypotheses,

\[
\boxed{\quad
|t|\,|\sigma+2t|
\le C(K,T)\sqrt{d_{\sigma,t}(1+B+d_{\sigma,t})}.
\quad}                                                     \tag{21}
\]

In particular, for \(d_{\sigma,t}\le1\),

\[
\operatorname{dist}(t,\{0,-\sigma/2\})^2
\le C(K,T)(1+B)d_{\sigma,t}.                               \tag{22}
\]

To verify this, replace the scalar target \(-4\) by \(4\sigma\) in (17).
Its calibrated algebraic value becomes \[
P_\sigma(t)=-\frac{16}{3}t^2(\sigma+2t).
\] The curvature budget uses \(|\operatorname{Ric}|\le|E_\sigma|+2\),
the amplitude estimate is unchanged, and weak Bianchi differentiates
neither constant multiple of the metric. Equations (16), (18) and (18a)
therefore hold with this \(P_\sigma\) and \(d_{\sigma,t}\); division by
\(|t|\) proves (21), including its trivial \(t=0\) case. The distance
estimate follows as before. For \(|t|\ge t_0>0\), the same
good/bad-amplitude argument controls
\(\fint(|\nabla\psi|^2+|d\psi|^2)\).

At zero defect and \(t\ne0\), this forces \(t=-\sigma/2\) and a parallel
norm-squared-two self-dual form. The metric is then Kähler--Einstein
with \(\operatorname{Ric}=\sigma g\). The positive Kähler--Einstein
amplitude is \(-1/2\), not \(+1/2\). This extension concerns the
analytic criterion only: it does not supply positive-side monopole
classes or a Yau26 positive-branch obstruction.

END COMPLETE SOURCE

\begin{center}\rule{0.5\linewidth}{0.5pt}\end{center}

\hypertarget{source-3-line1_functional_pencil.md}{%
\subsection{Source 3:
line1\_functional\_pencil.md}\label{source-3-line1_functional_pencil.md}}

Original path: mathematics/current/line1\_functional\_pencil.md

SHA-256:
13f0324dab4999ffba998cd60e8a1d8c152b97aae6f40218aca599230c4201ae

BEGIN COMPLETE SOURCE

\hypertarget{quantitative-non-saturation-of-scalarweyl-functionals-near-negative-einstein-metrics}{%
\section{Quantitative non-saturation of scalar--Weyl functionals near
negative Einstein
metrics}\label{quantitative-non-saturation-of-scalarweyl-functionals-near-negative-einstein-metrics}}

23 September 2026

\hypertarget{main-theorem}{%
\subsection{1. Main theorem}\label{main-theorem}}

Let \((M,g)\) be a closed connected oriented smooth four-manifold. Fix a
Spin\(^c\) structure, with real first Chern class \(a\), whose
unperturbed Seiberg--Witten equations have a solution for every metric
conformal to \(g\). The structure can depend on \(g\); it is fixed while
the conformal weight varies. Write \[
V=\operatorname{Vol}_g(M),\qquad p=(a_g^+)^2>0,\qquad
r=\fint_M|\operatorname{Ric}+g|^2,\qquad
b_a=\frac{(-a^2)_+}{V}.
\] Ricci uses the tensor norm, two-forms the exterior-algebra norm, and
\(W^+\) the Hilbert--Schmidt norm as an endomorphism of \(\Lambda^+\).
Set \(w=\sqrt6|W^+|\). For a constant \(\beta>0\), define the
dimensionless, nonnegative functional deficit \[
\mathfrak D_\beta(g,a)
=\frac{\beta+1}{\beta}\int_M(s^2+\beta w^2)
-72\pi^2p.
\tag{1}
\] Here dimensionless means invariant under constant metric rescaling;
\(\mathfrak D_\beta/V\), like \(r\), is measured after fixing the
Einstein target to \(-g\).

\textbf{Theorem 1 (functional non-saturation).} For every finite \(L>2\)
there is \(c_L>0\) such that, simultaneously for \(2\le\beta\le L\), \[
\boxed{
(1+b_a)\left(r+\frac{\mathfrak D_\beta(g,a)}V\right)
\ge c_L(\beta-2)^2.
}
\tag{2}
\] There is no scalar-curvature band, Sobolev constant,
injectivity-radius, diameter, volume, or full-curvature hypothesis.
Actual conformal monopole solvability and the displayed square cost are
the inputs.

If the selected class additionally satisfies
\(a^2\ge T:=2\chi(M)+3\tau(M)\), there is \(c'_L>0\) with \[
\boxed{
r+\frac{\mathfrak D_\beta(g,a)}V\ge c'_L(\beta-2)^2.
}
\tag{3}
\] In particular, for an exact metric \(\operatorname{Ric}=-g\) and such
a class, \[
\int(s^2+6\beta|W^+|^2)
\ge\frac{72\pi^2\beta}{\beta+1}p
+c'_L\frac{\beta}{\beta+1}(\beta-2)^2V.
\tag{4}
\]

The endpoint \(\beta=2\) must remain: a negative Kähler--Einstein
surface with its canonical Spin\(^c\) structure has \(r=0\) and
\(\mathfrak D_2=0\). Thus (2) is a compatibility theorem, not a
prohibition of the Kähler--Einstein branch.

The proof supplies all auxiliary forms and their joint analytic budget
from (1). It extends the former calibration at \(\beta=8\) to a
parameter family of curvature functionals. It does not, by itself,
identify a new smooth topology with no Einstein metric.

\hypertarget{the-exact-deficit-decomposition}{%
\subsection{2. The exact deficit
decomposition}\label{the-exact-deficit-decomposition}}

Put \[
q=-s+w,\quad q_+=\max(q,0),\quad q_-=\max(-q,0),\quad
\Delta_+=\int q_+^2-72\pi^2p.
\] The signed monopole estimate gives \(\Delta_+\ge0\). The relevant
classical inequality follows from the rescaled equations; the positive
part is needed when scalar curvature changes sign. Completing a square
gives the exact identity \[
\boxed{
\mathfrak D_\beta
=\Delta_++\int q_-^2+\frac1\beta\int(s+\beta w)^2.
}
\tag{5}
\] Indeed
\((\beta+1)(s^2+\beta w^2)/\beta =(-s+w)^2+(s+\beta w)^2/\beta\).

Write \(e=s+4\), \(z=s+\beta w\), and \[
q_\beta=4\frac{\beta+1}{\beta},\qquad
\kappa=r+\mathfrak D_\beta/V.
\] Then \[
w-4/\beta=(z-e)/\beta,\qquad
q-q_\beta=z/\beta-(1+1/\beta)e,\qquad
\int e^2\le4Vr.
\tag{6}
\] Consequently, uniformly for \(2\le\beta\le L\), \[
\int\bigl[(q-q_\beta)^2+(w-4/\beta)^2\bigr]
\le C_L V\kappa,
\quad
\int(|\operatorname{Ric}|^2+s^2+|W^+|^2)
\le C_L V(1+\kappa).
\tag{7}
\] These bounds control the Weyl norm, not its eigenspaces. The latter
come from the actual monopole below.

\hypertarget{actual-bounded-spinors-with-no-scalar-band}{%
\subsection{3. Actual bounded spinors with no scalar
band}\label{actual-bounded-spinors-with-no-scalar-band}}

Set \[
Q=\max(q,3),\qquad
f_0=\max\{\sqrt Q,(-s)_+/4\}.
\] Choose smooth positive majorants \(f_j\to f_0\) uniformly on the
fixed compact manifold. Conformal covariance, applied to
\(\widehat g=f_j^{-2}g\), supplies solutions \[
D_A\Phi=0,\quad -iF_A^+=f_j\sigma(\Phi),\quad
\psi=2\sqrt2\sigma(\Phi),\quad
\rho=|\Phi|^2=|\psi|,\quad u=3\rho/2.
\tag{8}
\] We use \(|\sigma(\Phi)|^2=|\Phi|^4/8\). The spinor equation is
\(4\nabla_A^*\nabla_A\Phi+(s+f_j\rho)\Phi=0\). At a maximum of \(\rho\),
\(f_j\rho\le-s\), so \[
0\le\rho\le4,\qquad 0\le u\le6.
\tag{9}
\] Every approximation error below tends to zero for the fixed \(g,a\).
No uniform smoothing rate along a sequence and no bounds on derivatives
of \(f_j\) are required.

The cutoff costs only scalar error: \[
0\le Q^2-q_+^2\le9e^2,\qquad
0\le f_0^2-Q\le\frac9{16}e^2.
\tag{10}
\] For the second inequality, a nonzero difference forces
\((-s)_+>4\sqrt3>6\), whence \((-s)_+\le3|e|\). For the first, \(q<3\)
implies \(s>-3\), hence \(e>1\).

Define nonnegative defects \[
\begin{aligned}
D_0={}&2\int|d\rho|^2
+\int(4\rho|\nabla_A\Phi|^2-|\nabla\psi|^2)\\
&+2\int|d\psi|^2+2\int E_\lambda,\qquad
E_\lambda=\langle W^+\psi,\psi\rangle+w\rho^2/3\ge0.
\end{aligned}
\tag{11}
\] Nonnegativity uses the differential inequality for the parallel
spinor quadratic map and \(\lambda_{\min}(W^+)\ge-w/3\). The spinor and
self-dual-form Weitzenböck identities, subtracted after testing the
spinor against \(\rho\Phi\), give \[
\frac{27}{8}D_0=\int(qu^2-f_j u^3).
\tag{12}
\]

For completeness, the moment argument that retains the truncation cost
is as follows. Set \[
C=\int Qu^2,\quad Z=\int\sqrt Q\,u^3,\quad E_Q=\int Q^2.
\] Then \(C-Z\ge27D_0/8\ge0\), while the curvature equation gives \[
C_{\rm act}:=\int f_j^2u^2=18\|F_A^+\|_2^2\ge72\pi^2p,
\quad 0\le C_{\rm act}-C\le\frac{81}{4}\int e^2+o(1).
\] Hölder gives \(C\le E_Q^{1/3}Z^{2/3}\), hence \[
0\le E_Q-C\le S+o(1),\qquad
S:=\Delta_++\frac{117}{4}\int e^2\le117V\kappa.
\] If \(\theta=(Z/E_Q)^{1/3}\), the nonnegative Hölder slack implies
\(E_Q-C\ge2(C-Z)\). The associated amplitude estimate is
\(\int Q(u-\sqrt Q)^2\le E_Q-C\): with probability measure
\(Q^2d\mu/E_Q\), let \(h=u/\sqrt Q\), \(c=\mathbb E h^2\),
\(z_3=\mathbb E h^3\), and \(m=\mathbb E h\). Then
\(z_3\le c\le z_3^{2/3}\), so \(c\le1\); Cauchy--Schwarz gives
\(c^2\le m z_3\le mc\), and hence \(m\ge c\). Thus
\(\mathbb E(h-1)^2\le1-c\). We have proved \[
D_0\le\frac{4S}{27}+o(1),\qquad
\int Q(u-\sqrt Q)^2\le S+o(1).
\tag{13}
\]

\hypertarget{produce-the-full-calibrated-form-including-its-weyl-direction}{%
\subsection{4. Produce the full calibrated form, including its Weyl
direction}\label{produce-the-full-calibrated-form-including-its-weyl-direction}}

Because \(Q\ge3\), \(u\le6\), and \((Q-q)^2\le e^2\), (13) yields \[
\int(q-u^2)^2
\le2(1+2\sqrt3)^2S+2\int e^2+o(1)
\le C V\kappa+o(1).
\tag{14}
\] The middle inequality uses
\(|Q-u^2|\le(1+2\sqrt3)\sqrt Q\,|\sqrt Q-u|\). Together with (7), this
controls \(\int(u^2-q_\beta)^2\).

Now choose the normalization, rather than assuming it: \[
c_\beta^2=\frac{9\beta}{8(\beta+1)},\qquad
\omega=c_\beta\psi,\qquad t_\beta=1/\beta,\qquad
\alpha_\beta=\frac9{8(\beta+1)}.
\] It follows that \[
\|\omega\|_\infty\le3\sqrt2,\qquad
\int(|\omega|^2-2)^2\le C_L V\kappa+o(1).
\tag{15}
\] Define the smooth polynomial model \[
A_{\omega,t}=t\left(\frac{|\omega|^2}{3}I-\omega\otimes\omega\right).
\] The exact identity for arbitrary \(\psi\), including its zeros, is \[
\left|W^+-\alpha_\beta
       \left(\frac{\rho^2}{3}I-\psi\otimes\psi\right)\right|^2
=\frac16(w-2\alpha_\beta\rho^2)^2
+2\alpha_\beta E_\lambda.
\tag{16}
\] This identity uses only trace-freeness of \(W^+\), not an eigenline
selection. The tensor on its left is exactly \(W^+-A_{\omega,1/\beta}\).
Also \(2\alpha_\beta\rho^2=u^2/(\beta+1)\), and \[
w-\frac{u^2}{\beta+1}
=\left(w-\frac4\beta\right)
-\frac{u^2-q_\beta}{\beta+1}.
\] Equations (7), (13)--(16) give, for the same smooth solution, \[
d_\beta:=\fint\left(
 |\operatorname{Ric}+g|^2+|W^+-A_{\omega,1/\beta}|^2
+(|\omega|^2-2)^2\right)
\le C_L\kappa+o(1).
\tag{17}
\] In particular the small Weyl-direction error has been proved from the
actual gauge equations. It is not a hidden geometric hypothesis.

\hypertarget{supply-the-second-order-budget}{%
\subsection{5. Supply the second-order
budget}\label{supply-the-second-order-budget}}

For a bounded section \(v\) of a metric bundle with compatible
connection curvature \(\mathscr F\), covariant integration by parts on a
closed four-manifold gives \[
\|\nabla v\|_4^2\le4\|v\|_\infty\|\nabla^2v\|_2,
\] \[
\|\nabla^2v\|_2^2
\le2\|\nabla^*\nabla v\|_2^2
+32\|v\|_\infty^2\|\operatorname{Ric}\|_2^2
+33\|v\|_\infty^2\|\mathscr F\|_2^2.
\tag{18}
\] The curvature convention is
\(|\mathscr F|^2=\sum_{ij} |\mathscr F_{ij}|_{\rm HS}^2\). These are the
inherited Hessian identities, not a Sobolev embedding. One first
interchanges Hessian slots and then commutes on \(\nabla v\); no
derivative of curvature is introduced.

Equations (7), (10) give \(\int f_j^2\le C_L V(1+\kappa)+o(1)\). The
determinant Chern--Weil identity and (9) imply \[
\|F_A\|_2^2=2\|F_A^+\|_2^2-4\pi^2a^2
\le4\int f_j^2+4\pi^2(-a^2)_+.
\] The curvature of the positive-spinor connection is bounded in norm by
the \(W^+,s,\operatorname{Ric}^{\circ},F_A\) blocks; \(W^-\) acts
trivially on this representation. Its squared integral is therefore at
most \(C_LV(1+\kappa+b_a)+o(1)\). The spinor equation and \(|\Phi|\le2\)
likewise bound \(\int|\nabla_A^*\nabla_A\Phi|^2\) by
\(C_LV(1+\kappa)+o(1)\).

Apply (18) to \(\Phi\), and twice differentiate the parallel quadratic
map \(\Phi\mapsto\omega\). The resulting terms are bounded by
\(C(|\Phi||\nabla_A^2\Phi|+|\nabla_A\Phi|^2)\). Finally use
\(\Delta_H\omega=\nabla^*\nabla\omega+(s/3)\omega-2W^+\omega\). This
proves \[
B_\omega:=\fint|\Delta_H\omega|^2
\le C_L(1+\kappa+b_a)+o(1),
\] \[
\boxed{(1+B_\omega)d_\beta
\le C_L\kappa(1+\kappa+b_a)+o(1).}
\tag{19}
\] All the forms are smooth actual solutions. None of these estimates
differentiates the smoothing weight.

\hypertarget{spectral-compatibility-and-completion-of-the-proof}{%
\subsection{6. Spectral compatibility and completion of the
proof}\label{spectral-compatibility-and-completion-of-the-proof}}

The bounded self-dual-form compatibility theorem states that, for
\(\|\omega\|_\infty\le K\) and bounded constant \(t\), \[
|t|\,|1-2t|
\le C\sqrt{d_t(1+B_\omega+d_t)},
\tag{20}
\] where \(d_t\) is as in (17). Its proof is recorded in full in
\href{round21_analytic.md}{the analytic criterion}, Sections 2, 3 and 6.
The mechanism relevant to its use here is worth making explicit. The
same Hessian estimate bounds second derivatives of \(A_{\omega,t}\). The
weak contracted Bianchi identity, tested against
\(J_t=\operatorname{div}W(A_{\omega,t})\), gives \[
\fint|J_t|^2\le C|t|\sqrt{d_t(1+B_\omega+d_t)}.
\] If \(v=|\omega|^2\), the polynomial pointwise contraction gives \[
t^2v(|\nabla\omega|^2+|d\omega|^2)
\le |J_t|^2+\tfrac43t^2|dv|^2,
\quad
\fint|dv|^2\le C\sqrt{d_t(1+B_\omega+d_t)}.
\] Testing the self-dual Hodge equation against \(v\omega\) gives an
algebraic term \(16t^2(1-2t)/3\); all its derivative terms are
controlled by the preceding two displays. The zeroth-order errors cost
at most \(Ct^2\sqrt{d_t}\). Thus \[
\left|\frac{16}{3}t^2(1-2t)\right|
\le7\fint|J_t|^2+C t^2\sqrt{d_t(1+B_\omega+d_t)}.
\] Division by \(|t|\), with the case \(t=0\) separate, proves (20). The
argument neither differentiates the Ricci error nor assumes
\(d\omega=0\).

Substitute the produced form, with \(t=1/\beta\), into (20). Using
(17)--(19), then letting its smoothing allowance tend to zero, \[
\frac{(\beta-2)^2}{\beta^4}
\le C_L\kappa(1+\kappa+b_a).
\tag{21}
\] Absorb \(\beta^4\le L^4\) into the constant. When \(\kappa\le1\), the
right side is bounded by \(2C_L(1+b_a)\kappa\). When \(\kappa>1\), (2)
is immediate after choosing \(c_L\le(L-2)^{-2}\). This proves (2).

For (3), Chern--Weil gives \[
4\pi^2T=\int\left(2|W^+|^2+s^2/24
-|\operatorname{Ric}^{\circ}|^2/2\right)
\ge-\tfrac12 Vr.
\] Hence \(a^2\ge T\) implies \(b_a\le r/(8\pi^2)\). For \(\kappa\le1\)
this is bounded by \(1/(8\pi^2)\); for \(\kappa>1\) use the same trivial
split. Equation (4) is (3) with \(r=0\), after multiplying by
\(\beta V/(\beta+1)\).

\hypertarget{the-allowed-endpoint-without-a-scalar-curvature-band}{%
\subsection{7. The allowed endpoint without a scalar-curvature
band}\label{the-allowed-endpoint-without-a-scalar-curvature-band}}

The same construction has a structural consequence even when the
parameter approaches the allowed endpoint. It is not obtained by
declaring that endpoint forbidden.

Let \(H\) be the harmonic self-dual representative of \(a_g^+\), with
convention \(\operatorname{proj}_{\mathcal H^+}(iF_A^+)=2\pi H\). Set \[
\zeta_\beta=-\frac{3\pi\beta}{\beta+1}H.
\] \textbf{Theorem 2 (actual harmonic alignment).} Under the hypotheses
of Theorem 1, if \(\kappa\le1\), the same smooth forms can be chosen
with arbitrarily small allowance \(\eta>0\) so that \[
\fint|\omega-\zeta_\beta|^2\le C_L\kappa+\eta,
\quad
\fint|\nabla\omega|^2
\le C_L\sqrt{\kappa(1+\kappa+b_a)}+\eta.
\tag{22}
\] The amplitude, polynomial Weyl fit and Hodge-Laplacian bounds remain
(15), (17) and (19). In particular, at \(\beta=2\), the harmonic target
is \(-2\pi H\), and \(\kappa=r+144\pi^2D_{\rm LB}/V\). This conclusion
needs neither a scalar band nor a scalar spectral gap.

\textbf{Proof.} Retain the actual curvature moment from Section 3. Since
\(C\le E_Q\), its Hodge remainder satisfies \[
18\|iF_A^+-2\pi H\|_2^2
=C_{\rm act}-72\pi^2p\le S+o(1).
\] The curvature equation therefore gives \[
\|f_j\psi+4\sqrt2\pi H\|_2^2\le4S/9+o(1).
\tag{23}
\] Furthermore \[
\int(f_j-\sqrt{q_\beta})^2
\le2\int(f_j-\sqrt Q)^2
+2\int(\sqrt Q-\sqrt{q_\beta})^2
\le C_LV\kappa+o(1).
\] Here \((f_0-\sqrt Q)^2\le f_0^2-Q\le9e^2/16\), and the square-root
map is uniformly Lipschitz on \([3,\infty)\). Use \(|\psi|\le4\) and
multiply (23) by \(c_\beta^2/q_\beta\). The identity
\(4\sqrt2\pi c_\beta/\sqrt{q_\beta} =3\pi\beta/(\beta+1)\) proves the
first estimate in (22).

For the second, the derivative conclusion of the bounded-form theorem on
\(1/L\le t\le1/2\) is
\(\fint|\nabla\omega|^2\le C_L\sqrt{d_t(1+B_\omega+d_t)}\) when
\(d_t\le1\). If \(\kappa\) is small, (17) ensures this. For the
remaining compact range of \(\kappa\le1\), the integrated Hodge equation
and the bound on \(B_\omega\) give
\(\fint|\nabla\omega|^2\le C_L\sqrt{1+b_a}\); enlarging the constant
recovers (22). All allowances are taken small simultaneously for the
same solution.

The harmonic form need not have a uniform maximum bound. Thus (22) does
not license replacing \(\omega\) by \(\zeta_\beta\) inside the quadratic
Weyl model. It does give a precise global spectral conclusion: if
\(\Pi_{>\lambda}\) is the spectral projection of the rough Laplacian on
self-dual forms, then for every \(\lambda>0\), \[
\fint|\Pi_{>\lambda}\zeta_\beta|^2
\le C_L\left[\kappa+
\frac{\sqrt{\kappa(1+\kappa+b_a)}}{\lambda}\right].
\tag{24}
\] Indeed, project \(\zeta_\beta-\omega\) and use
\(\|\Pi_{>\lambda}\omega\|_2^2 \le\lambda^{-1}\|\nabla\omega\|_2^2\),
then let the allowance tend to zero. This controls an actual specified
harmonic class; it does not assert global Kähler topology or exact
metric correction.

At \(\beta=2\), there is a stronger direct estimate, independent of
\(b_a\). In the notation of (11), the retained Hodge identity is \[
\int|\nabla\psi|^2
=2\int|d\psi|^2+2\int E_\lambda
-\tfrac13\int(s+2w)\rho^2.
\] Since \(\rho\le4\), \(\int(s+2w)^2\le2\mathfrak D_2\), and
\(c_2^2=3/4\), this gives simultaneously \[
\fint|\nabla\omega|^2\le G(\kappa)+\eta,
\quad G(\kappa)=13\kappa+4\sqrt{2\kappa},
\quad
\fint|d\omega|^2\le13\kappa/2+\eta.
\tag{25}
\] The preceding harmonic-alignment estimate can likewise be made
explicit: \[
\fint|\omega+2\pi H|^2\le49\kappa+\eta.
\tag{26}
\] For (26), (6) gives \(\fint(q-6)^2\le18\kappa\). The cutoff estimate
implies \(\fint(f_j-\sqrt6)^2\le9\kappa+o(1)\), and (23) then costs at
most \(4\cdot9\kappa+S/(9V)\le49\kappa\). Thus in (24) the endpoint
right side can be replaced by \(98\kappa+2G(\kappa)/\lambda\), with no
square-cost hypothesis. The
\href{line1_composed_application.md}{endpoint application} gives this
statement and its spectral consequences separately.

\hypertarget{what-is-strengthened-and-what-is-not}{%
\subsection{8. What is strengthened, and what is
not}\label{what-is-strengthened-and-what-is-not}}

At \(\beta=2\), \[
\mathfrak D_2=144\pi^2D_{\rm LB},\qquad
D_{\rm LB}=\frac1{4\pi^2}\int(s^2/24+|W^+|^2/2)-p/2.
\] The classical equality set is negative cscK, compatible with the
selected Spin\(^c\) structure, not necessarily Einstein. Equation (2)
does not turn its zero deficit into an unconditional Ricci estimate.

At \(\beta=8\), Chern--Weil gives \[
\mathfrak D_8
=108\pi^2(T-2p/3)+\frac{27}{2}\int|\operatorname{Ric}^{\circ}|^2.
\tag{27}
\] Thus the old topological calibration is one member of the family. For
other parameters, (2) controls a different scalar--Weyl functional, with
its own genuine angular calibration supplied in Section 4. This is a
functional extension of the method, not a second discovery of the old
\(1/3\) blow-up obstruction.

An exact negative Kähler--Einstein metric with canonical class has
\(s=-4,w=2,p=V/(2\pi^2)\). Direct substitution yields \[
\frac{\mathfrak D_\beta}{V}
=\frac{4(\beta-2)^2}{\beta}.
\tag{28}
\] The quadratic vanishing at the allowed endpoint is consequently the
correct order; the theorem does not claim an optimal coefficient.

For \(0<\beta\le2\), the classical scalar and scalar--Weyl bounds
already provide a stronger elementary bound. If \(x=\|s\|_2,y=\|w\|_2\),
then \(x\ge\sqrt{32\pi^2p}\) and \(x+y\ge\sqrt{72\pi^2p}\). Minimizing
\(x^2+\beta y^2\) on these constraints gives
\(\int(s^2+\beta w^2) \ge8\pi^2p(4+\beta)\). Thus this range does not
need the new compatibility argument. For \(\beta>2\), that scalar
constraint is inactive at the classical relaxed minimizer. The
additional obstruction in (2) is differential, through weak Bianchi and
the actual bounded form.

For the Yau26 negative input, \(r\le4\delta_-(g)^2\). Under
\(a^2\ge T\), (3) therefore gives the usable estimate \[
4\delta_-(g)^2+\mathfrak D_\beta/V
\ge c'_L(\beta-2)^2.
\tag{29}
\] To obtain a new purely topological application one must additionally
force \(\mathfrak D_\beta/V\) small from independently justified
topological or geometric data. No such new topology is asserted here.
Neither (29) nor the allowed-root form supply proves existence of an
Einstein metric on an arbitrary smooth manifold.

\hypertarget{sources-and-verification}{%
\subsection{Sources and verification}\label{sources-and-verification}}

\begin{itemize}
\tightlist
\item
  C. LeBrun, \emph{Ricci Curvature, Minimal Volumes, and Seiberg--Witten
  Theory}, \href{https://arxiv.org/abs/math/0003068}{Theorem 2.3 and
  Propositions 3.1, 4.1}: the classical curvature inequalities and cscK
  equality set.
\item
  C. LeBrun, \emph{Four-manifolds, curvature bounds, and convex
  geometry}, \href{https://arxiv.org/abs/math/0611450}{Lemma 3.7}:
  weighted equations.
\item
  \href{round4_rigidity.md}{Integral truncated-weight proof}, Sections
  2--4: the inherited cutoff and moment estimates, reproduced above.
\item
  \href{round21_analytic.md}{Bounded-form spectral theorem}: the
  analytic compatibility theorem, including its polynomial contractions.
\item
  \href{line1_scalar_band_remainder.md}{Scalar-band supply}: the prior
  allowed cscK branch, which is not itself a forbidden-spectrum
  application.
\end{itemize}

Exact Einstein saturation forcing Kähler--Einstein structure already
appears in LeBrun's Corollary 2.5; it is not a new classification here.
The quantitative persistence under integral Ricci error, and the actual
all-scalar form-production estimates, are the assertions at issue.

The parameterized all-scalar functional estimate (2) is the claim proved
in this note. Its finite algebra is checked separately; those checks do
not replace the gauge-solvability hypothesis or the analytic proofs. No
comprehensive priority claim over later literature is made.

END COMPLETE SOURCE

\begin{center}\rule{0.5\linewidth}{0.5pt}\end{center}

\hypertarget{source-4-line1_compactified_pencil.md}{%
\subsection{Source 4:
line1\_compactified\_pencil.md}\label{source-4-line1_compactified_pencil.md}}

Original path: mathematics/current/line1\_compactified\_pencil.md

SHA-256:
0d971e1519c9d8b5edddcbd71f158c3ca08acf943d5ae4a7eb3f359af2628bed

BEGIN COMPLETE SOURCE

\hypertarget{a-compactified-functional-pencil-and-a-two-sided-einstein-energy-window}{%
\section{A compactified functional pencil and a two-sided Einstein
energy
window}\label{a-compactified-functional-pencil-and-a-two-sided-einstein-energy-window}}

23 September 2026. Baseline: \texttt{4401d28}.

This note extends the finite-parameter non-saturation theorem to its
vanishing-Weyl endpoint, with an explicit degenerating weight. It then
obtains the upper, as well as the lower, quantitative energy envelope.
The analytic input is the existing bounded-form compatibility theorem;
the classical feasible energy interval is not a new curvature
inequality.

\hypertarget{uniformity-on-the-unbounded-parameter-interval}{%
\subsection{1. Uniformity on the unbounded parameter
interval}\label{uniformity-on-the-unbounded-parameter-interval}}

Let \((M,g)\) be closed, connected, oriented and four-dimensional. Fix
an actual Spin\(^c\) structure whose ordinary Seiberg--Witten equations
are solvable for every metric conformal to \(g\). Write \[
 p=(a_g^+)^2>0,\quad V=\operatorname{Vol}(g),\quad
 r=\fint|\operatorname{Ric}+g|^2,\quad b_a=(-a^2)_+/V,
\] \[
 D_\beta=\frac{\beta+1}{\beta}\int(s^2+6\beta|W^+|^2)
                 -72\pi^2p,\qquad \beta\ge2.
\] The conventions are tensor Ricci norm and Hilbert--Schmidt Weyl
endomorphism norm. The class is fixed while the conformal weight varies.

\textbf{Theorem 1.} There is a universal \(c>0\) such that, for every
finite \(\beta\ge2\), \[
 \boxed{(1+b_a)(r+D_\beta/V)
       \ge c\frac{(\beta-2)^2}{\beta^4}.}                 \tag{1}
\] If \(a^2\ge T=2\chi+3\tau\), the factor \(1+b_a\) can be removed,
after reducing the universal constant. Unlike the earlier coefficient
\(c_L\), the constant in (1) does not depend on a parameter cutoff. Its
explicit weight vanishes at both the Kähler endpoint and the
vanishing-Weyl endpoint.

\textbf{Proof.} Put \(w=\sqrt6|W^+|\), \(q=-s+w\), \(e_s=s+4\),
\(z=s+\beta w\), and \(\kappa=r+D_\beta/V\). The exact nonnegative
decomposition in \href{line1_functional_pencil.md}{the finite-parameter
proof}, Section 2, implies \[
 \fint e_s^2\le4r,\quad \fint z^2\le\beta D_\beta/V,
 \quad\Delta_+/V\le D_\beta/V.
\] Consequently \[
 \fint(w-4/\beta)^2
 \le2D_\beta/(\beta V)+8r/\beta^2\le3\kappa,
\] \[
 \fint(q-q_\beta)^2\le14\kappa,
 \qquad q_\beta=4+4/\beta\in[4,6].                       \tag{2}
\] In particular the normalized Ricci, scalar and \(W^+\) squared
energies are bounded by \(C(1+\kappa)\), uniformly for \(\beta\ge2\).

Here is the required tracking of the form-production constants;
uniformity is not inferred merely by sending \(L\) to infinity. The
truncated weight \(Q=\max(q,3)\), \(f_0=\max(\sqrt Q,(-s)_+/4)\) is
independent of \(\beta\). The inherited moment argument gives the same
bounds \(|\psi|\le4\), \(D_0\le4S/27+o(1)\) and the same amplitude
control, with \(S\le117V\kappa\). Its errors involve only \(e_s\),
\(\Delta_+\) and \(Q\). For the normalized form and Weyl model the
coefficients are \[
 c_\beta^2=\frac{9\beta}{8(\beta+1)}\in[3/4,9/8],\quad
 \alpha_\beta=\frac9{8(\beta+1)}\le3/8,\quad
 t=1/\beta\in(0,1/2].
\] Set \(\omega=c_\beta\psi\) and
\(A_{\omega,t}=t(|\omega|^2I/3-\omega\otimes\omega)\). The polynomial
fit identity is \[
 |W^+-A_{\omega,t}|^2
 =\tfrac16(w-2\alpha_\beta|\psi|^2)^2
                                  +2\alpha_\beta E_\lambda.
\] The amplitude estimate, (2), and
\(2\alpha_\beta|\psi|^2=u^2/(\beta+1)\) therefore yield, with one smooth
actual solution and an arbitrarily small allowance, \[
 \|\omega\|_\infty\le3\sqrt2,\quad
 d_t:=\fint\bigl(|\operatorname{Ric}+g|^2
   +|W^+-A_{\omega,t}|^2+(|\omega|^2-2)^2\bigr)
                       \le C\kappa+o(1).                 \tag{3}
\] No inverse power of \(t\) is used here. The determinant energy
identity and the bounded-spinor Hessian estimate then give \[
 B_\omega:=\fint|\Delta_H\omega|^2
                   \le C(1+\kappa+b_a)+o(1),             \tag{4}
\] again uniformly: the needed curvature energies are bounded by (2),
and the normalization of the quadratic spinor map is bounded above. This
is exactly the estimate proved in Sections 3--5 of the finite-parameter
paper, with every parameter-dependent coefficient now bounded
explicitly.

The bounded-form compatibility theorem, with the fixed parameter range
\(0\le t\le1/2\), gives \[
 t^2(1-2t)^2\le C\kappa(1+\kappa+b_a).                   \tag{5}
\] For \(\kappa\le1\) this proves (1); for \(\kappa>1\) use
\(t^2(1-2t)^2\le1/64\). Let the smoothing allowance tend to zero. If
\(a^2\ge T\), Chern--Weil gives \(b_a\le r/(8\pi^2)\), and the same
split removes the square-cost factor. This proves the theorem. No
derivative estimate for the harmonic projection is asserted uniformly as
\(t\to0\). \(\square\)

\hypertarget{a-pencil-with-an-actual-endpoint-at-zero}{%
\subsection{2. A pencil with an actual endpoint at
zero}\label{a-pencil-with-an-actual-endpoint-at-zero}}

Define on \(0\le t\le1/2\) \[
 P_t=t(1+t)\int s^2+6(1+t)\int|W^+|^2-72\pi^2tp.
                                                               \tag{6}
\] For \(t>0\), this is \(tD_{1/t}\); at zero it is exactly
\(P_0=6\int|W^+|^2\). Hence \(P_t\ge0\) on the closed interval.
Multiplying (1) by \(t\), and using \(r+tD/V\ge t(r+D/V)\), proves \[
 \boxed{(1+b_a)(r+P_t/V)\ge c t^3(1-2t)^2.}             \tag{7}
\] The same square-cost removal applies when \(a^2\ge T\). At \(t=0\)
the inequality is trivial and supplies no nonzero lower amplitude or
actual anti-self-dual Einstein existence. At \(t=1/2\), the endpoint
includes negative Kähler--Einstein metrics. The cubic decay in (7) is a
proved sufficient rate from this transfer, not an optimality claim.

For a normalized negative Kähler--Einstein metric and its canonical
class, \[
 P_t/V=4(1-2t)^2.
\] This provides an actual normalization check. It does not establish
sharpness at the other endpoint.

\hypertarget{the-missing-upper-energy-envelope}{%
\subsection{3. The missing upper energy
envelope}\label{the-missing-upper-energy-envelope}}

Assume now that \(A>0\), \[
 T=2\chi+3\tau=A(1-h),\qquad 1/4<h<1/3,                 \tag{8}
\] and for every metric one can choose an actual conformally solvable
class with \(p\ge A\) and \(a^2\ge T\). Put \[
 e=\frac{\int|\operatorname{Ric}|^2}{8\pi^2A},\quad
 F(h)=\frac{(4-\sqrt{1-3h})^2}{9},\quad
 G(h)=\frac{(4+\sqrt{1-3h})^2}{9},
\] \[
 t_*(h)=\frac{1-2\sqrt{1-3h}}{8+2\sqrt{1-3h}}\in(0,1/8),
 \qquad \Theta(h)=t_*^3(1-2t_*)^2.                       \tag{9}
\] Use the all-sign relative defect \[
 \mathcal E(g)=\inf_{\lambda\ne0}
 \frac{\int|\operatorname{Ric}-\lambda g|^2}{\lambda^2V}.
\]

\textbf{Theorem 2.} There is a universal \(c_*>0\) such that every
smooth metric under (8) satisfies \[
 \boxed{\mathcal E(g)+(G(h)-e)_+\ge c_*\Theta(h).}        \tag{10}
\] Combining with the inherited lower-envelope theorem, every Einstein
candidate must satisfy the strictly smaller energy window \[
 \boxed{F(h)+\varepsilon h^2\le e
                       \le G(h)-c_*\Theta(h).}          \tag{11}
\] Ricci-flat and positive Einstein metrics are independently excluded
by the positive actual monopole scalar estimate. Equation (11) concerns
all remaining Einstein metrics, not just minimizers or a gluing
neighborhood.

\textbf{Proof.} Write \(R=\int|\operatorname{Ric}|^2\),
\(P=\int|\operatorname{Ric}^{\circ}|^2\), \(C_A=8\pi^2A\). The exact
Chern--Weil reduction of (6) is \[
 P_t=k(t)R+\ell(t)P+12\pi^2(1+t)T-72\pi^2tp,
\] \[
 k(t)=\frac{(1+t)(8t-1)}2,\qquad
 \ell(t)=2(1+t)(1-2t).                                  \tag{12}
\] For \(t=t_*(h)\), direct substitution gives \[
 k(t)G(h)+\tfrac32(1+t)(1-h)-9t=0,
\] and therefore \[
 P_t=C_A k(t)(e-G(h))+\ell(t)P-72\pi^2t(p-A)
       \le C_A|k(t)|(G(h)-e)_++\ell(t)P.                \tag{13}
\] This is the opposite sign of the coefficient in the lower-envelope
argument. It is why the finite-parameter range beyond \(\beta=8\)
controls an upper energy wall.

First fix the target to \(\operatorname{Ric}=-g\) and
\(r=\fint|\operatorname{Ric}+g|^2\le1\). The scalar monopole bound and
the trace estimate imply \[
 C_A/V\le(2+\sqrt r)^2\le9,\qquad P/V\le r.
\] On \(0\le t\le1/8\), \(|k(t)|\le1/2\) and \(0\le\ell(t)\le2\). Using
(7), with its square-cost factor removed, in (13) gives \[
 c\Theta(h)\le3r+\tfrac92(G(h)-e)_+
              \le\tfrac92\{r+(G(h)-e)_+\}.             \tag{14}
\] For \(r>1\) the desired estimate follows by decreasing the constant.
Constant rescaling and taking the infimum over negative target constants
prove (10) for the negative-target infimum \(\mathcal E_-\).

The passage to all signs needs a different argument from the earlier
lower-envelope transfer: a very large \(e\) now makes the energy excess
zero. Here the upper Chern--Weil budget closes that possibility. Let
\(I=\int s\). If \(I\le0\), direct minimization gives
\(\mathcal E_+=4\). If \(I>0\), it gives
\(\mathcal E_+=\mathcal E=4-I^2/(RV)\). The signed scalar estimate
\(\int s_-^2\ge4C_A\) implies \[
 \mathcal E_+\ge4/e.                                   \tag{15}
\] If \(\mathcal E_+<1/2\), the chirality upper bound gives \[
 e(1-\mathcal E_+)\le
 U=3\min(2\chi+3\tau,2\chi-3\tau)/A\le3,
\] so \(e\le6\). Equation (15) would imply \(\mathcal E_+\ge2/3\), a
contradiction. Thus \(\mathcal E_+\ge1/2\). Since \(\Theta\) is bounded,
this proves (10) also for the positive target infimum, after reducing
\(c_*\); take the minimum of the two signwise bounds. \(\square\)

\hypertarget{geometry-of-the-two-endpoints-and-an-actual-application}{%
\subsection{4. Geometry of the two endpoints and an actual
application}\label{geometry-of-the-two-endpoints-and-an-actual-application}}

The classical Einstein constraints are, with nonnegative \(x,y\), \[
 x^2=4e,\quad x^2+8y^2=12(1-h),\quad x\ge2,\quad x+y\ge3.
\] When \(1/4<h<1/3\), the ellipse has \(x<3\). Solving these
inequalities gives the full classical interval \(F(h)\le e\le G(h)\).
Thus the interval itself is classical; the positive quantitative margins
in (11) are the new application of the differential non-saturation
mechanism.

At \(h\downarrow1/4\), the upper optimal parameter tends to \(t=0\), and
\(\Theta(h)\sim(8/27)(h-1/4)^3\). At \(h\uparrow1/3\), both optimal
parameters meet \(t=1/8\); the interval collapses, consistently with the
old critical gap. At the lower-envelope endpoint \(h\downarrow0\), the
optimal parameter instead tends to the allowed Kähler value \(t=1/2\).
These are different geometric degeneracies; they cannot be treated as
one compact parameter range with a fixed forbidden distance.

For an actual example, let \[
 M=(\Sigma_2\times\Sigma_3)\#5\overline{\mathbb{CP}}^{\,2}.
\] Its polarized actual classes give \(A=16\), \(h=5/16\), \(a^2=T=11\).
Here \[
 F=25/16,\qquad G=289/144,\qquad t_*=1/17,
 \qquad \Theta=225/17^5.
\] Any Einstein metric, if one exists, must consequently satisfy \[
 \frac{25}{16}+\varepsilon\frac{25}{256}
 \le\frac{R}{128\pi^2}
 \le\frac{289}{144}-c_*\frac{225}{17^5}.                 \tag{16}
\] This is not a nonexistence claim for this example. Its bare
two-chirality upper bound is only \(e\le33/16\), so the upper envelope
contains information not supplied by that bound alone. The actual
monopole selection is the same established polarized input as in the
main paper; no new gauge existence theorem is asserted here.

\hypertarget{relation-to-the-research-goal}{%
\subsection{5. Relation to the research
goal}\label{relation-to-the-research-goal}}

The first-line theorem now supplies a two-sided restriction on the
entire energy spectrum of possible Einstein outputs in the displayed
regime. It does not merely study an energy-minimizing sequence. For a
pointwise input, \(\mathcal E\le4\delta_\sigma^2\) converts (10) into
the corresponding upper-energy constraint on that input; a small
residual alone does not force the input energy into a forbidden region.

This note neither classifies the allowed spectra of all Einstein metrics
nor constructs an Einstein metric. The separate Bauer--Furuta
application in this increment is the new topological-category result.
The present contribution is the closed-parameter formulation and its
upper-wall application, using the existing analytic mechanism.

Sources: \href{https://arxiv.org/abs/math/0003068}{LeBrun's scalar--Weyl
inequalities and Chern--Weil reduction}; the local
\href{line1_functional_pencil.md}{functional-pencil proof}, Sections
2--6; \href{round21_analytic.md}{bounded-form compatibility}, Section 6;
\href{line1_energy_envelope_gap.md}{lower energy envelope}; and
\href{line1_two_chirality_topology.md}{two-chirality upper bound}.

\texttt{check\_line1\_compactified\_pencil.py} verifies the exact
coefficients, both envelope roots, endpoint rates, actual example, and
several deliberate sign errors. These checks do not substitute for the
written uniform analytic estimates or actual gauge-solvability inputs.

END COMPLETE SOURCE

\begin{center}\rule{0.5\linewidth}{0.5pt}\end{center}

\hypertarget{source-5-line1_energy_envelope_gap.md}{%
\subsection{Source 5:
line1\_energy\_envelope\_gap.md}\label{source-5-line1_energy_envelope_gap.md}}

Original path: mathematics/current/line1\_energy\_envelope\_gap.md

SHA-256:
8cd0c429b2f403b1c31bb24333efc7397afba5742394a2951dd56aa5e24c3ad4

BEGIN COMPLETE SOURCE

\hypertarget{a-uniform-ricci-energy-gap-above-the-classical-einstein-envelope}{%
\section{A uniform Ricci-energy gap above the classical Einstein
envelope}\label{a-uniform-ricci-energy-gap-above-the-classical-einstein-envelope}}

23 September 2026. This is an application of the parameterized
functional theorem, with actual polarized monopole classes. The
elementary energy envelope below follows from classical inequalities.
The added conclusion is a uniform positive separation from that
envelope, persisting under integral Einstein error. No new topological
nonexistence theorem is asserted.

\hypertarget{statement-on-a-realized-family}{%
\subsection{1. Statement on a realized
family}\label{statement-on-a-realized-family}}

Let \[
 M=X\#k\overline{\mathbb {CP}}^{\,2}\#m(S^1\times S^3),\qquad
 A=c_1^2(X)>0,\quad h=\frac{k+4m}{A},\quad 0<h<\frac13,
 \tag{1}
\] where \(X\) is a minimal complex surface of general type. Put \[
 R=\int_M|\operatorname{Ric}|^2,\qquad
 e=\frac{R}{8\pi^2A},\qquad
 F(h)=\frac{(4-\sqrt{1-3h})^2}{9},
\] \[
 \mathcal E(g)=\inf_{\lambda\ne0}
 \frac{\int_M|\operatorname{Ric}-\lambda g|^2}
      {\lambda^2\operatorname{Vol}(g)}.
 \tag{2}
\] Both \(e\) and \(\mathcal E\) are invariant under constant rescaling.

\textbf{Theorem.} Relative to Theorem 1 of
\href{line1_functional_pencil.md}{the functional-pencil proof}, there is
a universal \(\varepsilon>0\), independent of \(X,k,m,h,g\), such that
\[
 \boxed{\mathcal E(g)+(e-F(h))_+\ge\varepsilon h^2.}
 \tag{3}
\] Thus, for example, \(\mathcal E(g)\le\varepsilon h^2/2\) implies
\(e\ge F(h)+\varepsilon h^2/2\). Every Einstein metric on \(M\), if one
exists, necessarily satisfies \[
 \boxed{\frac{R}{8\pi^2A}\ge F(h)+\varepsilon h^2.}
 \tag{4}
\] An actual nonzero monopole projection excludes Einstein constants
zero and positive, so (4) concerns all Einstein candidates, not just a
chosen sign.

The same theorem holds whenever \(T=2\chi+3\tau=A(1-h)\) and each metric
admits an actual conformally solvable Spin\(^c\) structure with
\(p=(a_g^+)^2\ge A\) and \(a^2\ge T\). No regularity of the class
selection in \(g\) is needed. In (1), the polarized sign selection gives
\[
 p\ge A,\qquad a^2=A-k\ge A-k-4m=T.
 \tag{5}
\] This is gauge-theoretic input, not just an intersection-form
calculation; the conformal class keeps its polarization and its selected
structure. For \(b_+=1\), the small-perturbation chamber and, for
handles, the loop insertions are those in the
\href{lebrun_stability_closure.md}{polarized closure}. The source is
\href{https://arxiv.org/pdf/math/0003068}{LeBrun, Proposition 3.2 and
its proof}.

For a concrete family without chamber issues, take
\(X=\Sigma_g\times\Sigma_q\), \(g,q\ge2\). Then \(A=8(g-1)(q-1)\) and
\(b_+>1\). In particular,
\((\Sigma_2\times\Sigma_2)\#2\overline{\mathbb {CP}}^{\,2}\) has
\(A=8,h=1/4\). Formula (3) applies to every smooth metric on this actual
closed manifold. Any Einstein metric on it would satisfy \[
 \frac{R}{64\pi^2}\ge\frac{49}{36}+\frac{\varepsilon}{16}.
\] We do not assume or produce an Einstein metric on it.

\hypertarget{the-classical-envelope-is-not-the-new-result}{%
\subsection{2. The classical envelope is not the new
result}\label{the-classical-envelope-is-not-the-new-result}}

For an Einstein metric, set \[
 x^2=\frac{\int s^2}{8\pi^2A},\qquad
 y^2=\frac{6\int|W^+|^2}{8\pi^2A}.
\] The classical scalar and scalar--Weyl inequalities and Chern--Weil
give \[
 x\ge2,\qquad x+y\ge3,\qquad x^2+8y^2=12(1-h),\qquad e=x^2/4.
 \tag{6}
\] These imply \(x\ge(8-2\sqrt{1-3h})/3\), hence \(e\ge F(h)\). Equality
is feasible in this finite relaxation; it is not claimed to be realized
by a metric. The constraints are the scalar estimate,
\href{https://arxiv.org/pdf/math/0003068}{LeBrun's Theorem 2.3}, and the
exact Chern--Weil identity. This optimization is entirely classical.

The all-metrics infimum in (1) is \(\inf e=1+h\), by
\href{https://arxiv.org/pdf/math/0003068}{LeBrun's Theorem 4.3}. One has
the exact comparison \[
 F(h)-(1+h)=\frac{4h^2}{(1+\sqrt{1-3h})^2}>0.
 \tag{7}
\] The new term in (4) is therefore an additional uniform energy
separation above the already optimized classical restriction, not the
separation in (7). Classical equality rigidity already excludes exact
Einstein saturation away from the Kähler--Einstein endpoint; it does not
alone give a uniform gap for varying manifolds or almost-Einstein
metrics.

\hypertarget{exact-chernweil-reduction-and-the-proof}{%
\subsection{3. Exact Chern--Weil reduction and the
proof}\label{exact-chernweil-reduction-and-the-proof}}

Use tensor Ricci norms and Hilbert--Schmidt Weyl operator norms. Write
\[
 P=\int|\operatorname{Ric}^{\circ}|^2,\quad S_2=\int s^2,\quad
 W_2=\int|W^+|^2.
\] The identities \[
 S_2=4(R-P),\qquad W_2=2\pi^2T-R/12+P/3
\] turn the functional in the companion theorem into \[
 \mathfrak D_\beta
 =K_\beta R+L_\beta P+12\pi^2(\beta+1)T-72\pi^2p,
 \tag{8}
\] where \[
 K_\beta=\frac{(\beta+1)(8-\beta)}{2\beta},\qquad
 L_\beta=\frac{2(\beta+1)(\beta-2)}{\beta}.
\] Let \(z=\sqrt{1-3h}\) and choose \[
 \beta=\beta(h)=\frac{8-2z}{1+2z}\in(2,8).
 \tag{9}
\] Direct substitution gives \[
 K_\beta F(h)+\tfrac32(\beta+1)(1-h)-9=0.
\] Using (5) in (8), with no discarded negative error, yields \[
 \mathfrak D_\beta
 =8\pi^2A K_\beta(e-F(h))+L_\beta P-72\pi^2(p-A)
 \le8\pi^2A K_\beta(e-F(h))+L_\beta P.
 \tag{10}
\]

First fix the Einstein target to \(-g\), and let
\(V=\operatorname{Vol}(g)\), \(r=V^{-1}\int|\operatorname{Ric}+g|^2\).
The functional-pencil theorem, with \(L=8\) and \(a^2\ge T\), supplies a
universal \(c>0\) such that \[
 c(\beta-2)^2\le r+\mathfrak D_\beta/V.
\] Since \(P/V\le r\), equation (10) implies \[
 c(\beta-2)^2
 \le (1+L_\beta)r+
 K_\beta\frac{8\pi^2A}{V}(e-F(h)).
 \tag{11}
\] The scalar monopole bound and the trace inequality give \[
 32\pi^2A\le S_2,\qquad
 \sqrt{S_2/V}\le4+2\sqrt r,
 \qquad \frac{8\pi^2A}{V}\le(2+\sqrt r)^2.
 \tag{12}
\] For \(r\le1\), replace the last factor in (11) by its positive part
before using (12). On \([2,8]\), \[
 0\le K_\beta\le9/2,\qquad 0\le L_\beta\le27/2.
\] Consequently \[
 c(\beta-2)^2\le\frac{81}{2}\bigl[r+(e-F(h))_+\bigr].
 \tag{13}
\] Finally, \[
 \beta(h)-2
 =\frac{18h}{(1+z)(1+2z)}\ge3h.
\] Thus (3), with \(r\) in place of \(\mathcal E\), holds for \(r\le1\)
with coefficient \(2c/9\). For \(r>1\) it is immediate after choosing
\(\varepsilon=\min\{2c/9,1\}\), because \(h^2<1/9\). For any
\(\lambda<0\), apply this result to \((-\lambda)g\): its normalized
error is the corresponding quotient in (2), while \(e\) and \(h\) are
unchanged. Taking the infimum proves (3) first for \(\mathcal E_-\), the
infimum restricted to negative \(\lambda\).

To include positive target constants, put \(C=8\pi^2A\), \(I=\int s\),
and let \(\mathcal E_+\) be the infimum for \(\lambda>0\). The signed
scalar estimate is \[
 \int s_-^2\ge32\pi^2p\ge4C,\qquad s_-:=\max\{-s,0\}.
 \tag{14}
\] For clarity, it follows directly from an ordinary monopole solution:
with \(\rho=|\Phi|^2\), the integrated spinor identity gives
\(\int\rho^2\le\int s_-\rho\), while Hodge projection gives
\(\int\rho^2\ge32\pi^2p\); apply Cauchy--Schwarz. If \(I>0\), then \[
 \frac{I^2}{V}\le\int s_+^2\le S_2-4C\le4R-4C,
 \qquad \mathcal E_+=4-\frac{I^2}{RV}\ge\frac4e.
\] If \(I\le0\), direct minimization gives \(\mathcal E_+=4\). When
\(e\le F(h)+1\le25/9\), therefore, \(\mathcal E_+\ge36/25>1\). When
\(e>F(h)+1\), the energy excess alone exceeds one. Hence
\(\mathcal E_++(e-F(h))_+\ge1\ge\varepsilon h^2\). Together with
\(\mathcal E=\min\{\mathcal E_-,\mathcal E_+\}\), this proves (3) for
all target signs.

Only the functional-pencil theorem is the nonclassical analytic input.
There is no geometric compactness, scalar band, auxiliary-form
assumption, or additional square-cost budget in this transfer.

\hypertarget{why-the-parameter-eight-does-not-give-this-energy-separation}{%
\subsection{4. Why the parameter eight does not give this energy
separation}\label{why-the-parameter-eight-does-not-give-this-energy-separation}}

At \(\beta=8\), \(K_8=0\). Equation (8) loses \(R\) altogether: \[
 \mathfrak D_8=108\pi^2(T-2p/3)+27P/2.
 \tag{15}
\] This is the old topological gap. Varying \(\beta\) retains the
independent energy coordinate and gives the separation in (3).

There is an exact algebraic non-implication check. At \(h=1/4\), set
\(p=A\), \(a^2=T=3A/4\), \(P=0\), \(s=-4\), and prescribe the constant
Weyl norm \(w=\sqrt6|W^+|=8/7\). Choose \[
 V=2\pi^2A(49/36),\qquad R=4V.
\] These moment data satisfy Chern--Weil, the scalar inequality, every
classical bound \(\mathfrak D_\gamma\ge0\), \(\gamma>0\), and
\(e=F(1/4)=49/36\). However, \[
 \mathfrak D_{7/2}=0,\qquad \mathfrak D_8/V=162/49.
 \tag{16}
\] They violate the new parameter-\(7/2\) estimate. They do not violate
any possible universal coefficient in the global parameter-eight member
\(r+\mathfrak D_8/V\ge36c_8\): necessarily \(c_8\le1/11\). Indeed, for
an actual negative Kähler--Einstein metric \(g_0\), the rescalings
\(g=tg_0\) give \[
 r+\mathfrak D_8/V=4(1-1/t)^2+18/t^2
 \ge36/11,
\] with equality at \(t=11/2\). Yet \(162/49>36/11\). This proves
independence from the parameter-eight numerical inequality even at the
largest coefficient it could possibly have.

The data in (16) are a formal curvature-moment model, not a metric
construction. The manifold and its actual monopole-class selection in
Section 1 are realized; realizing the forbidden moments would contradict
the theorem. This argument does not claim independence from every
possible use of the old differential proof machinery.

\hypertarget{the-remaining-topological-input-is-precise}{%
\subsection{5. The remaining topological input is
precise}\label{the-remaining-topological-input-is-precise}}

Equation (3) does not exclude all energies on any new smooth topology.
For normalized Einstein metrics, the missing independent number is \[
 e=\frac{R}{8\pi^2A}=\frac{V}{2\pi^2A}.
 \tag{17}
\] It is not fixed by \(T\), \(A\), or \(p\). Chern--Weil fixes a linear
combination of scalar and Weyl energies, not this scalar/volume
coordinate. An independently proved upper bound
\(e<F(h)+\varepsilon h^2\) for every Einstein candidate would convert
(4) into a new topological obstruction; no such bound is supplied here.
An all-metrics energy \emph{infimum} cannot serve as that upper bound.

For the pointwise negative Ricci input
\(\delta_-=\|\operatorname{Ric}^{\sharp}+I\|_{\infty,\mathrm{op}}\), (3)
implies the raw consequence \[
 4\delta_-^2+(e-F(h))_+\ge\varepsilon h^2.
 \tag{18}
\] Thus a construction with controlled energy below the strengthened
envelope cannot be arbitrarily Ricci-pinched. This is a restriction on
such input metrics, not a proof that their smooth manifold has no
Einstein metric at a different energy.

\texttt{check\_line1\_energy\_envelope\_gap.py} checks the exact
identities, envelope optimization, constants, scale transformation, and
formal non-implication witness. These finite checks do not certify
Seiberg--Witten solvability or the inherited differential theorem. No
priority claim over unexamined later literature is made.

END COMPLETE SOURCE

\end{document}
